\documentclass[11pt]{article}

\usepackage[T1]{fontenc}
\usepackage[utf8]{inputenc}
\usepackage[margin=1in]{geometry}
\usepackage{lmodern}
\usepackage{microtype}
\usepackage{amsmath,amssymb,mathtools,amsthm}
\usepackage{booktabs,tabularx,array,longtable,multirow}
\usepackage[section]{placeins}
\usepackage{float}
\usepackage{graphicx}
\usepackage{xcolor}
\usepackage{enumitem}
\usepackage{needspace}
\usepackage{url}
\usepackage[numbers,sort&compress]{natbib}
\usepackage[hidelinks]{hyperref}
\usepackage[nameinlink,noabbrev]{cleveref}

\allowdisplaybreaks
\setlength{\emergencystretch}{2em}
\raggedbottom
\setlist{leftmargin=*,topsep=3pt,itemsep=1.5pt,parsep=0pt,partopsep=0pt}

\definecolor{deepblue}{RGB}{23,69,124}
\hypersetup{
  colorlinks=true,
  linkcolor=deepblue,
  citecolor=deepblue,
  urlcolor=deepblue,
  pdftitle={A Symmetry Quotient Attack on Odd-Characteristic SNOVA},
  pdfauthor={Justin Thaler, Kai Zhang, Scott Kominers, Quang Dao},
  pdfsubject={Cryptanalysis of odd-characteristic SNOVA and its Version 2.3 parameters},
  pdfkeywords={SNOVA, multivariate signatures, symmetric public matrices, forgery, NIST PQC}
}

\newtheorem{theorem}{Theorem}[section]
\newtheorem{lemma}[theorem]{Lemma}
\newtheorem{proposition}[theorem]{Proposition}
\newtheorem{corollary}[theorem]{Corollary}
\newtheorem{definition}[theorem]{Definition}
\newtheorem{remark}[theorem]{Remark}

\newcommand{\F}{\mathbb F}
\newcommand{\Mat}{\operatorname{Mat}}
\newcommand{\Sym}{\operatorname{Sym}}
\newcommand{\rank}{\operatorname{rank}}
\newcommand{\Span}{\operatorname{span}}
\newcommand{\im}{\operatorname{im}}
\newcommand{\transpose}{\mathsf T}
\newcommand{\SNOVA}{\textsf{SNOVA}}
\newcommand{\AXN}{\mathsf{AXN}}
\newcommand{\attacktitle}{A Symmetry Quotient Attack on Odd-Characteristic
  \texorpdfstring{\SNOVA}{SNOVA}}

\title{\attacktitle}
\author{%
  Justin Thaler\thanks{a16z crypto research and Georgetown University} \and
  Kai Zhang\thanks{MIT} \and
  Scott Kominers\thanks{Harvard University and a16z crypto research} \and
  Quang Dao\thanks{Carnegie Mellon University and LayerZero Labs}%
}
\date{22 August 2026}

\begin{document}
\maketitle

\begin{abstract}
\SNOVA{} compresses UOV public keys by representing variables and coefficients
with small matrices.  Its current odd-characteristic design uses symmetric
public quadratic matrices.  We show that this symmetry creates a new weakness
in the common-column restriction used by affine-column forgery attacks.

If the public matrix algebra has dimension $d$, the restricted verifier has
only $d(d+1)/2$ distinct quadratic labels rather than $d^2$: sixteen forms
collapse to ten for $\ell=4$, and four to three for $\ell=2$.  An invertible
public output transformation separates these surviving quadratics from the
affine constraints and target checks.  For every published Version~2.3
parameter shape passing the stated public rank checks, this gives an exact
square quadratic system while preserving every verifier equation, rejection
rule, and final verification check.

For the $\ell=2$ Level-I shape we give a separate Boolean attack.  On one
official KAT key, a deterministic $(15,4,7)$ MOR staircase and an exact
rank-$24$ diagonal-pencil certificate are supplied; the same staircase
succeeds on all $100$ archived Level-I KAT public keys.  Target-averaged
counting handles every quadratic branch and singular residual system without
May--Ostuzzi--Ressler Assumption~1.  After conditioning on the complete
diagonal transcript, the remaining Hermitian channel yields a full-rank joint
polar pencil except with conditional probability below $2^{-114.684}$ in the
pinned XOF-transcript model.  On every such completion, the expected
success-adjusted work is below $2^{132.447}$ Boolean gates, $10.553$ bits below
the Level-I classical reference, with explicitly bounded low working memory.

For the other parameter rows, the recovery estimates remain conditional on a
symbolic-homotopy hypothesis and an explicit unresolved multiplier
$\kappa_{\rm hom}$.  We do not claim a fresh-key probability for the diagonal
rank certificate or a production wall-clock forgery.
\end{abstract}


\section{Introduction}\label{sec:intro}

NIST's Additional Digital Signatures process seeks signature schemes with
security and performance profiles that complement the standards selected in
its first post-quantum process.  \SNOVA{} is one of the nine candidates
advanced to the third round in May 2026~\cite{NISTRound3,NISTIR8610}.  It
belongs to the Unbalanced Oil and Vinegar line of multivariate signatures,
where verification evaluates a public quadratic map and signing uses a
hidden central map with an oil-and-vinegar structure
~\cite{KipnisPatarinGoubin1999,SNOVA23}.

\SNOVA{} obtains compact public keys by representing variables and
coefficients with small matrices and generating many public forms from a
smaller structured family.  That structure has led to key-recovery and direct
forgery attacks.  The former expose the hidden UOV structure, while Beullens
uses a public relation among signature columns.  Later work improves both
lines through group actions, extension-field lifting, and
multihomogeneous solving
~\cite{IkematsuAkiyama2024,LiDing2024,Beullens2024SNOVA,Cabarcas2025SNOVA,FurueEtAl2025,NakamuraEtAl2025}.
These works analyze the original or Round~2 construction over $\F_{16}$.

The next revision answered Ran's wedge attack, which uses the alternating
polar forms available in characteristic two.  Bros et al. adapt the method to
SNOVA's block structure~\cite{Ran2026,BrosEtAl2026SNOVAWedge}.  The \SNOVA{}
team responded with odd-prime fields and symmetric public quadratic matrices,
while warning that symmetrization could create new attacks
~\cite{WangEtAl2025Round2Security}.  NIST later described this change as a
response to the wedge attack and did not regard \SNOVA{} as stable
~\cite{NISTIR8610}.

This change entered the specification in Version~2.1, not Version~2.3.
Version~2.3 subsequently added rectangular signatures and revised parameter
sets~\cite{SNOVA23}.  Our quotient theorem concerns the symmetric
odd-characteristic construction and is not specific to those nine parameter
sets.  The concrete rank, root, and cost analysis covers the nine published
$q=19$ Version~2.3 shapes.

Prior attacks already cover odd characteristic.  Jin et al. give a UOV
key-recovery attack in arbitrary characteristic~\cite{JinEtAl2026}.  Ran gives
an odd-characteristic construction using symmetric forms, but SNOVA's added
structure requires separate analysis~\cite{Ran2026}.  Ding, Li, and Tao give a
wedge framework valid in any characteristic, while their concrete SNOVA
estimates cover the even-characteristic Round~2 parameters
~\cite{DingLiTao2026}.  None combines the symmetric public matrices in the
current $q=19$ construction with the affine-column forgery relation.

This brings us to the question addressed in this paper.  What happens when an
affine-column attacker imposes the same common-column relation on the new
symmetric verifier?  We prove that symmetry makes pairs of restricted
quadratic forms identical.  An invertible public change of output coordinates
then separates these quadratics from the remaining affine constraints and
target checks.  For every published $q=19$ parameter shape that passes the
stated public rank checks, this leaves a quadratic system with the same number
of equations and variables.

This statement concerns the restricted verifier, not the unrestricted
matrix-valued map.  Let $A=\F_q[S]$ have dimension $d$, and let
$\Sigma_\xi=I_n\otimes\xi$ for $\xi\in A$.  After all signature columns are
expressed through one stacked column $u$, symmetry gives
$q_{i,\xi,\eta}(u)=q_{i,\eta,\xi}(u)$.  The $d^2$ ordered labels therefore
factor through $\Sym^2(A)$, with the exact dimension change
\begin{equation}
 d^2\ \longrightarrow\ \binom{d+1}{2},
 \qquad 16\ \longrightarrow\ 10\ \text{ for }d=4,
 \qquad 4\ \longrightarrow\ 3\ \text{ for }d=2.
 \label{eq:intro-duplication}
\end{equation}
The full verifier relates the two ordered features by matrix transposition;
only the common-column restriction makes them the same scalar polynomial.
The identification summarized in \Cref{eq:intro-duplication} is exact and
basis-independent, but deliberately local to the attacker-imposed
restriction (\Cref{thm:symmetry-quotient}).

\subsection{Our results}

\paragraph{An exact symmetry quotient and public reduction.}
After the common-column restriction, every symmetric public base form factors
through $\Sym^2(A)$.  This loss occurs before the outer public mixing, so later
linear combinations cannot restore the missing forms.  Public row reduction
then separates the surviving quadratics from the affine conditions and target
checks.  For every published parameter shape that passes the stated public
rank checks, this gives an exact square quadratic system.  Every recovered
root is reconstructed and checked by the original verifier
(\Cref{thm:symmetry-quotient,prop:exact-affine,cor:exact-square-restriction}).

\paragraph{Recovery bounds, with assumptions separated.}
The exact reduction does not by itself prove that an acceptable root exists or
that a square system can be solved at the advertised cost.  For the homotopy
route we therefore keep the unproved hypothesis $\mathsf H_{\rm hom}$ explicit
and report arithmetic ledgers before the unresolved multiplier
$\kappa_{\rm hom}$.  \Cref{thm:sq-all-nine} gives the corresponding row-by-row
break-even values rather than assigning that multiplier a guessed constant.

The Level-I $(48,16,2,2)$ row has a stronger, separate analysis.  For one
official key we supply a target-independent $(15,4,7)$ MOR staircase and an
exact rank-$24$ diagonal-pencil certificate.  Aggregate target averaging then
gives expected trial work below $2^{131.447}$ without
May--Ostuzzi--Ressler Assumption~1.  Conditional on the diagonal transcript,
\Cref{thm:jg-joint-fullrank} bounds failure of the joint $D/H$ full-rank
predicate by $2^{-114.684}$ in the pinned Hermitian-transcript model.  On every
completion satisfying that predicate, the exact joint second moment gives
success probability $>1/2$ and expected success-adjusted work below
$2^{132.447}$.  The deterministic staircase/rank preflight costs below
$2^{50}$ AXN gates and the online enumeration is depth first with explicitly
bounded small mutable state.  We do not infer a fresh-key pass probability for
the diagonal rank certificate from the finite KAT evidence.

For the Level-III and Level-V $\ell=2$ rows, the homotopy ledgers remain below
their respective NIST references provided
$\log_2\kappa_{\rm hom}<23.205$ and $38.156$.  No homotopy exponent in this
paper is a measured wall-clock or peak-memory bound.

\paragraph{Design implications.}
We derive necessary output-dimension conditions for redesigns that retain the
common-column architecture.  These conditions are not sufficient for security
(\Cref{thm:upgrade-combined-floor,thm:upgrade-zero-offset-floor}).  The dimension floors identify where the exact quotient persists, independently
of whether the unresolved solver multipliers are small enough to cross a NIST
cost target.

\subsection{How to read the result}

The reduction, root-existence argument, and solver bound have different
logical status.  We use the following convention throughout.

\begin{table}[H]
\centering
\caption{Logical status of the main claims.}
\label{tab:logical-status}
\footnotesize
\renewcommand{\arraystretch}{1.04}
\begin{tabularx}{0.98\textwidth}{@{}>{\raggedright\arraybackslash}p{0.24\textwidth}>{\raggedright\arraybackslash}p{0.18\textwidth}X@{}}
\toprule
Part & Status & Meaning\\
\midrule
Symmetry quotient and output split & Exact after public rank checks & Deterministic public algebra; every verifier equation and rejection check is preserved.\\
Level-I $\ell=2$ Just-Guess route & Exact diagonal certificate + conditional Hermitian theorem & One official key has the staircase and rank-$24$ certificate.  Conditioned on them, the bad joint-pencil event has probability $<2^{-114.684}$ and every good completion costs $<2^{132.447}$ in expectation.\\
Other recovery rows & Conditional & Homotopy completeness assumes $\mathsf H_{\rm hom}$ and retains the unresolved multiplier $\kappa_{\rm hom}$.\\
Implementation claims & Split & The Level-I Just-Guess traversal is explicitly low-space; no production wall-clock solve is claimed, and homotopy peak memory remains unresolved.\\
\bottomrule
\end{tabularx}
\end{table}

Every returned candidate is reconstructed and checked by the original verifier.
The deterministic $(15,4,7)$ staircase also succeeds on all $100$ archived
Level-I KAT public keys, but we do not infer a fresh-key probability for the
rank-$24$ certificate from that finite census.

\subsection{Attack overview}

In the hidden UOV coordinates, \SNOVA{} contains no product of two oil
blocks.  The signer first chooses the vinegar blocks, after which the signing
equations are linear in the oil entries.  A secret invertible change of
variables hides this division, so the public verifier presents a structured
quadratic map rather than the linear system used by the signer
(\Cref{sec:scheme}).

The attack seeks a direct public preimage, not the hidden oil space.  It starts
from the common-column relation used in earlier forgery attacks.  Inside every
matrix-valued signature block, the attacker writes all columns as public
affine functions of one stacked column.  This selects a public subset of
possible signatures; it is not an identity of the unrestricted verifier
(\Cref{sec:public-symmetry}).

After this restriction, every ordered pair of public matrix multipliers gives
a scalar quadratic polynomial.  Symmetry makes the polynomials for the two
orders equal.  One group of sixteen ordered pairs therefore contributes at
most ten distinct quadratics when $\ell=4$, while one group of four
contributes at most three when $\ell=2$.  The equality holds before the outer
public map combines these forms, so later linear combinations cannot restore
the missing quadratic directions (\Cref{thm:symmetry-quotient}).

The attacker next computes an invertible public change of verifier outputs.
Some new coordinates contain the surviving quadratic forms.  The remaining
coordinates are affine conditions on the signature and, when necessary,
consistency conditions on the target.  Solving the affine equations gives a
public affine space.  Restricting this space to as many variables as surviving
quadratics gives a square system without dropping any verifier equation
(\Cref{prop:exact-affine,cor:exact-square-restriction}).

There are then two recovery choices.  The complete system solves every
residual quadratic equation on a square slice.  A smaller system solves only
selected equations on a lower-dimensional slice and checks every omitted
equation afterward.  The smaller system has a much smaller algebraic state,
but a random trial hits the full target less often.  An exact public
determinant test permits the smaller system; if that test fails, the current
algorithm uses the complete system (\Cref{sec:sq-adaptive}).

The root analysis bounds the chance that a trial contains an accepted root
with a full-rank Jacobian.  The solver analysis is separate.  It starts from a
polynomial system with known roots and symbolically deforms that system into
the \SNOVA{} residual system.  The hypothesis $\mathsf H_{\rm hom}$ requires
the deformation and its specialization to behave regularly enough for the
solver to recover the nonsingular roots within the stated arithmetic cost.
This hypothesis affects completeness and cost, but not the validity of a
returned signature (\Cref{app:conditional-solver}).

Finally, every recovered root is converted back into a candidate signature,
checked against the public format rule, and passed to the original verifier.
A failed rank test, target check, solve, or final verification causes a
restart.  The final check makes the attack one-sided: it can fail to return a
signature, but it never returns an invalid one.

\section{The \SNOVA{} signature scheme}\label{sec:scheme}

This section defines the part of Version~2.3 used in the attack.  Fix
parameters $(v,o,q,\ell,r)$ and put $n=v+o$.  A vector of private or public
variables has the form
\[
 X=(X_1,\ldots,X_n)\in\Mat_{\ell\times r}(\F_q)^n.
\]
In the private coordinates, $X_1,\ldots,X_v$ are the vinegar blocks and
$X_{v+1},\ldots,X_n$ are the oil blocks.  The central map has $o$ outputs in
$\Mat_{r\times\ell}(\F_q)$, so both its oil input and its output contain
$or\ell$ field elements~\cite{KipnisPatarinGoubin1999,SNOVA23}.

\paragraph{The hidden central map.}
The UOV structure excludes every product of two oil variables.  Write
$[n]=\{1,\ldots,n\}$.  The allowed ordered pairs are
\begin{equation}
 \Omega=[n]^2\setminus\{v+1,\ldots,n\}^2.
 \label{eq:scheme-omega}
\end{equation}
Let $S\in\Mat_\ell(\F_q)$ be the public symmetric matrix with irreducible
characteristic polynomial, and let
\[
 A=\F_q[S]\cong\F_{q^\ell}.
\]
Every element of $A$ is a symmetric matrix, and all elements of $A$ commute.
Version~2.3 uses
\[
 m_1=\left\lceil\frac{or}{\ell}\right\rceil,
 \qquad N_\alpha=r\ell+r
\]
base forms and outer summands.  For each $h\in[m_1]$, the secret base form is
an $n$ by $n$ block matrix $[F_h]=(F_{h,jk})$, where
$F_{h,jk}\in\Mat_\ell(\F_q)$ and $F_{h,jk}=0$ whenever
$(j,k)\notin\Omega$.

We write $L_{i,\alpha}$ and $R_{i,\alpha}$ for the matrices called
$A_{i,\alpha}$ and $B_{i,\alpha}$ in the specification.  Their dimensions are
$r$ by $r$ and $r$ by $\ell$, respectively.  The other public factors
$Q_{i,\alpha,1},Q_{i,\alpha,2}$ are nonzero elements of $A$.  Put
$h(i,\alpha)=1+((i-1+\alpha)\bmod m_1)$.  The $i$th component of the hidden
central map is
\begin{equation}
 \widetilde F_i(X)=
 \sum_{\alpha=0}^{N_\alpha-1} L_{i,\alpha}
 \left(
   \sum_{(j,k)\in\Omega}
   X_j^\transpose
   Q_{i,\alpha,1}F_{h(i,\alpha),jk}Q_{i,\alpha,2}X_k
 \right)R_{i,\alpha},
 \qquad i\in[o].
 \label{eq:scheme-central}
\end{equation}
The inner sum is an $r$ by $r$ matrix, and the full output is an $r$ by
$\ell$ matrix.  Since \Cref{eq:scheme-central} has no product of two oil
blocks, fixing the vinegar blocks makes it linear in every entry of the oil
blocks.

\paragraph{The secret change of variables and the public map.}
The secret matrix $T^{12}=(T_{ab})$ has $v$ by $o$ entries in $A$.  It defines
the invertible change of variables $X=T(U)$ by
\begin{equation}
 \begin{aligned}
  X_a&=U_a-\sum_{b=1}^o T_{ab}U_{v+b}, &&1\le a\le v,\\
  X_{v+b}&=U_{v+b}, &&1\le b\le o.
 \end{aligned}
 \label{eq:scheme-secret-change}
\end{equation}
Thus $T$ hides the private vinegar and oil coordinates while preserving the
matrix shape of each block.  For every base form, key generation computes
\begin{equation}
 [P_h]=[T]^\transpose[F_h][T].
 \label{eq:scheme-base-public}
\end{equation}
The commutativity of $A$ then gives the public quadratic map
\begin{equation}
 \begin{split}
 \widetilde P_i(U)
 &=\sum_{\alpha=0}^{N_\alpha-1} L_{i,\alpha}
 \left(
   \sum_{j,k=1}^n
   U_j^\transpose
   Q_{i,\alpha,1}P_{h(i,\alpha),jk}Q_{i,\alpha,2}U_k
 \right)R_{i,\alpha}\\
 &=\widetilde F_i(T(U)).
 \end{split}
 \label{eq:scheme-public-map}
\end{equation}
For odd $q$, the public base forms are symmetric.  In block notation, this
means
\begin{equation}
 P_{h,kj}=P_{h,jk}^\transpose.
 \label{eq:scheme-public-symmetry}
\end{equation}
Equivalently, expanding all block entries over $\F_q$ gives a symmetric
$n\ell$ by $n\ell$ matrix $P_h$.  This is the symmetry used by our attack.

The block equations also explain the public key compression.  Partition the
base forms according to the vinegar and oil coordinates:
\[
 [F_h]=
 \begin{bmatrix}F_h^{11}&F_h^{12}\\F_h^{21}&0\end{bmatrix},
 \qquad
 [T]=
 \begin{bmatrix}I&-T^{12}\\0&I\end{bmatrix},
 \qquad
 [P_h]=
 \begin{bmatrix}P_h^{11}&P_h^{12}\\P_h^{21}&P_h^{22}\end{bmatrix}.
\]
The public seed determines the outer matrices, the $Q$ matrices, and
$P_h^{11},P_h^{12},P_h^{21}$.  The private seed determines $T^{12}$.  Key
generation reconstructs the hidden blocks and the remaining public block by
\begin{equation}
 \begin{aligned}
  F_h^{11}&=P_h^{11},
  &F_h^{12}&=P_h^{12}+P_h^{11}T^{12},\\
  F_h^{21}&=P_h^{21}+(T^{12})^\transpose P_h^{11},
  &P_h^{22}&=-(T^{12})^\transpose F_h^{12}-P_h^{21}T^{12}.
 \end{aligned}
 \label{eq:scheme-keygen-blocks}
\end{equation}
Thus the compressed public key is the public seed together with the blocks
$P_h^{22}$ for $h\in[m_1]$.  The secret key contains the seeds needed to
regenerate $T$ and the hidden forms~\cite{SNOVA23}.

\paragraph{Signing and verification.}
For a message $M$ and a 16-byte salt, the target is
\begin{equation}
 Y=H\bigl(s_{\rm public}\mathbin\|\operatorname{SHAKE256}(M,64)
          \mathbin\|\mathrm{salt}\bigr)
 \in\Mat_{r\times\ell}(\F_q)^o.
 \label{eq:scheme-target}
\end{equation}
The signer derives the vinegar blocks from the private seed, $M$, the salt,
and an attempt counter.  Write $V=(X_1,\ldots,X_v)$ for the chosen vinegar
blocks and let $z$ list the $or\ell$ entries of the oil blocks.  By
\Cref{eq:scheme-omega}, the signing equation is a square linear system
\begin{equation}
 G_Vz=\operatorname{vec}(Y)-c_V,
 \label{eq:scheme-signing-system}
\end{equation}
where $c_V=\operatorname{vec}(\widetilde F(V,0))$ and $G_V$ is the coefficient
matrix of the terms containing one vinegar block and one oil block.  If $G_V$
is singular, the signer chooses a new attempt.  Otherwise it solves for the
oil blocks, sets $U=T^{-1}(X)$, and retries if $U$ violates the public format
rule.  The signature is $U\mathbin\|\mathrm{salt}$.

The verifier expands the public key, checks the same format rule, recomputes
$Y$, and accepts exactly when
\begin{equation}
 \widetilde P(U)=Y.
 \label{eq:scheme-verification}
\end{equation}
Version~2.1 introduced the odd-prime fields and
\Cref{eq:scheme-public-symmetry}.  Version~2.3 retained that symmetry and
allowed $r\ne\ell$, which gives rectangular signature blocks.

\paragraph{Column form of the public verifier.}
For each column $c\in\{0,\ldots,r-1\}$, stack the $c$th columns of the public
signature blocks as
\begin{equation}
 u_c=\bigl(U_1[:,c]^\transpose,\ldots,U_n[:,c]^\transpose\bigr)^\transpose
 \in\F_q^{n\ell}.
 \label{eq:scheme-stacked-columns}
\end{equation}
Using the scalar expansion of $[P_h]$ and
$\Sigma_\xi=I_n\otimes\xi$, each scalar entry inside
\Cref{eq:scheme-public-map} is a linear combination of terms
\begin{equation}
 u_c^\transpose\Sigma_\xi P_h\Sigma_\eta u_d,
 \qquad \xi,\eta\in A.
 \label{eq:scheme-power-pair}
\end{equation}
The outer matrices $L_{i,\alpha}$ and $R_{i,\alpha}$ combine these terms by a
public linear map.  Thus SNOVA's matrix structure affects the scalar verifier
through the ordered left and right algebra multipliers in
\Cref{eq:scheme-power-pair}.  Section~4 imposes the attacker's column relation
and determines which of these ordered terms remain distinct.

\section{Prior attacks and mathematical background}\label{sec:background}

\subsection{SNOVA's attack and revision history}

The first line of \SNOVA{} cryptanalysis attacks the hidden UOV structure.
Ikematsu and Akiyama analyze key recovery on the scalar expansion of the core
map, and Li and Ding sharpen the interpretation as a UOV system with
$\ell v$ vinegar variables and $\ell o$ oil variables over the base field
~\cite{IkematsuAkiyama2024,LiDing2024}.  Nakamura, Tani, and Furue instead
lift the same structure to a smaller UOV instance over $\F_{q^\ell}$
~\cite{NakamuraEtAl2025}.  These attacks seek the hidden oil space or an
equivalent signing key.  They concern the original characteristic-two
construction and do not use the symmetric public matrices introduced later.

Beullens exposes a different layer.  He rewrites \SNOVA{} as a structured
MAYO-style whipping of a bilinear map, chooses an affine relation among the
signature columns, and searches for rank-deficient combinations in the outer
coefficient map~\cite{Beullens2024SNOVA}.  The result is a forgery attack,
including weak-key instances much cheaper than the generic case.  The Round~2
changes vary the outer coefficient matrices and mitigate the largest rank
drops, but they retain the affine-column attack surface.  Our attack starts
from this public restriction and acts one layer earlier: symmetry identifies
the restricted power-pair polynomials before the outer map sees them.

Cabarcas et al. and Furue et al. improve the algebraic solving of the systems
arising in the earlier attacks~\cite{Cabarcas2025SNOVA,FurueEtAl2025}.  The
former exploits stability under a commutative group action and adapts the
result to reconciliation and direct forgery.  The latter writes the
extension-field transformation explicitly and applies the resulting
multihomogeneous systems to intersection and rectangular-MinRank key recovery.
Their estimate for the old Round~2 $(q,v,o,\ell)=(16,29,6,5)$ set is below the
claimed Level-V threshold, subject to their stated solving-degree heuristic.
Both papers study the characteristic-two construction.  Neither derives a
loss of independent verifier equations from the later symmetric public
matrices.

Ran's wedge attack begins from the alternating polar forms of
characteristic-two UOV and recovers the hidden oil space through exterior
algebra~\cite{Ran2026}.  His latest version also gives a different construction
for odd characteristic, but states that SNOVA's added structure prevents the
analysis from carrying over directly.  Bros et al. specialize the
characteristic-two method to \SNOVA{} and give conjectural estimates below the
claimed security levels for every Round~2 parameter set
~\cite{BrosEtAl2026SNOVAWedge}.  The \SNOVA{} team's later analysis replaces
the conjectural kernel formula with a generating function and finds six
Round~2 sets below target~\cite{TsengEtAl2026Wedge}.

Ding, Li, and Tao place these attacks in a common Macaulay-matrix framework and
extend it to multihomogeneous systems in any characteristic
~\cite{DingLiTao2026}.  Their formulation applies to SNOVA's hidden oil-space
problem, but their concrete SNOVA analysis is expressly limited to the
even-characteristic Round~2 parameters.  It does not analyze the symmetric
public matrices or the nine $q=19$ Version~2.3 shapes.

Jin et al. give the main prior attack that genuinely crosses the
characteristic boundary~\cite{JinEtAl2026}.  Their $p$-reduced symmetric
algebra recovers Ran's exterior algebra at $p=2$ and defines an XL
key-recovery attack for arbitrary characteristic.  Their concrete
odd-characteristic application is to the twelve QR-UOV parameter sets, where
the attack does not beat reconciliation on the recommended parameters.  The
method targets the hidden UOV oil space.  It does not analyze the symmetric
public-matrix encoding of odd-characteristic \SNOVA{}, and it does not produce
the verifier-preserving direct-preimage reduction used here.

The distinction between these attacks is summarized in
\Cref{tab:prior-attacks}.  The exact new statement in this paper is the
$\Sym^2(A)$ factorization after the public common-column restriction and the
resulting invertible change of verifier outputs.  The direct-forgery cost
estimates built on that reduction remain conditional on their stated
root-count and solver hypotheses.

\begin{table}[H]
\centering
\caption{Scope of the closest attacks.  ``Output'' is what a successful attack
recovers; it is not a claim that every cited estimate is practical.}
\label{tab:prior-attacks}
\footnotesize
\renewcommand{\arraystretch}{1.08}
\begin{tabularx}{0.98\textwidth}{@{}>{\raggedright\arraybackslash}p{0.19\textwidth}>{\raggedright\arraybackslash}p{0.16\textwidth}>{\raggedright\arraybackslash}p{0.25\textwidth}>{\raggedright\arraybackslash}X@{}}
\toprule
Work & Targeted regime & Main mechanism & Output and relation to this paper\\
\midrule
Ikematsu--Akiyama; Li--Ding; Nakamura--Tani--Furue
& Original / early $q=16$
& Scalar or extension-field UOV interpretation
& Oil-space or equivalent-key recovery; no odd-$q$ public symmetrization.\\
Beullens; Cabarcas et al.; Furue et al.
& Original and Round~2 $q=16$
& Affine-column rank drop, group action, and multihomogeneous solving
& Direct forgery or key recovery; supplies the public restriction that this paper composes with symmetry.\\
Ran; Bros et al.; Tseng et al.; Ding--Li--Tao
& Arbitrary-characteristic UOV frameworks; concrete SNOVA estimates for Round~2 $q=16$
& Exterior or symmetric forms, Macaulay matrices, and the wedge-map kernel
& Oil-space recovery; no concrete analysis of the symmetric $q=19$ SNOVA verifier.\\
Jin et al.
& UOV families in arbitrary characteristic; concrete odd-$q$ QR-UOV instances
& $p$-reduced symmetric-algebra XL
& Oil-space recovery; no analysis of symmetric odd-$q$ SNOVA public matrices.\\
This paper
& All nine $q=19$ Version~2.3 shapes
& Common-column restriction, exact $\Sym^2(A)$ quotient, and public output split
& Conditional direct preimages checked by the original verifier; no hidden-key recovery claim.\\
\bottomrule
\end{tabularx}
\end{table}

The complementary generic-MQ and partitioning literature includes
Hashimoto's underdefined optimizer, Just Guess, generalized partitioning, and
recent F4 engineering~\cite{Hashimoto2023,MayOstuzziRessler2026,AsanumaEtAl2026,SakataTakagi2026}.

\subsection{MQ systems, roots, and affine slices}
An MQ system over $\F_q$ is a map $g:\F_q^N\to\F_q^K$ whose coordinates
have total degree at most two, together with a target $z\in\F_q^K$.  Linear
and constant terms are allowed.  A root is an $x$ satisfying $g(x)=z$.
Underdetermination makes roots more plentiful but does not make the system
linear or automatically easy.

An affine space is a translate $a+V$ of a linear subspace.  Restricting the
search to an affine slice $a+L\subseteq a+V$ turns the problem into a system
in $h=\dim L$ coordinates.  Three events must be separated:
\begin{enumerate}[label=(\roman*)]
\item the target--slice pair contains a base-field root;
\item it contains an accepted root at which a chosen Jacobian has full rank;
\item the chosen algebraic solver recovers every such root within its stated
      work bound.
\end{enumerate}
The first two are properties of the map and slice.  The third is a solver
statement.  The later recovery analysis proves them separately.

\section{From the public verifier to a square system}
\label{sec:public-symmetry}\label{sec:affine-reduction}

This section starts from the column form of the public verifier in
\Cref{eq:scheme-stacked-columns,eq:scheme-power-pair} and gives the exact attack
as one sequence of public operations.  The attacker imposes the common-column
relation.  Symmetry then identifies restricted quadratic forms.  An invertible
change of verifier outputs separates the remaining quadratic equations from
the affine ones.  Once two public rank checks pass, a final affine restriction
leaves equally many quadratic equations and variables.  This section makes no
claim that the resulting system has a root or that a solver can find one.

\subsection{Parameters and public forms}
Retain the notation of Section~2 and put
\begin{equation}
 m_1=\left\lceil\frac{or}{\ell}\right\rceil,
 \qquad M=or\ell,
 \qquad N_0=n\ell,
 \qquad K=m_1\binom{\ell+1}{2}.                    \label{eq:parameters}
\end{equation}
The ceiling follows the pinned Version~2.3 reference
implementation~\cite{SNOVA23}.  The theorem below proves that $K$ is the
maximum number of distinct restricted quadratic forms left by symmetry.

For the public rows, $A=\F_q[S]\cong\F_{q^\ell}$ and
$\dim_{\F_q}A=\ell$.  Write $\Sigma_\xi=I_n\otimes\xi$ for $\xi\in A$ and
use the symmetric scalar expansions
$P_1,\ldots,P_{m_1}\in\Mat_{N_0}(\F_q)$ from
\Cref{eq:scheme-base-public}.  For $\xi,\eta\in A$, define
\begin{equation}
  B_i^{\xi,\eta}(x,y)=x^\transpose\Sigma_\xi P_i\Sigma_\eta y.   \label{eq:base-bilinear}
\end{equation}
Setting $x=y$ gives the quadratic power-pair feature.

\begin{table}[H]
\centering
\caption{Dimensions of the nine public $q=19$ Version~2.3 parameter rows.}
\label{tab:parameters}
\small
\begin{tabular}{@{}c l r r r r@{}}
\toprule
Level & $(v,o,\ell,r)$ & $m_1$ & $M$ & $K$ & $N_0$\\
\midrule
I   & $(28,5,4,4)$  & $5$  & $80$  & $50$ & $132$\\
I   & $(48,16,2,2)$ & $16$ & $64$  & $48$ & $128$\\
I   & $(28,4,4,5)$  & $5$  & $80$  & $50$ & $128$\\
III & $(40,7,4,4)$  & $7$  & $112$ & $70$ & $188$\\
III & $(72,24,2,2)$ & $24$ & $96$  & $72$ & $192$\\
III & $(38,5,4,5)$  & $7$  & $100$ & $70$ & $172$\\
V   & $(50,9,4,4)$  & $9$  & $144$ & $90$ & $236$\\
V   & $(96,32,2,2)$ & $32$ & $128$ & $96$ & $256$\\
V   & $(52,6,4,6)$  & $9$  & $144$ & $90$ & $232$\\
\bottomrule
\end{tabular}
\end{table}

\subsection{The common-column restriction}
Write the stacked signature columns as
\[
  U=[u_0\mid\cdots\mid u_{r-1}],\qquad u_j\in\F_q^{N_0}.
\]
The attacker chooses public multipliers $R_j\in A$ and offsets
$v_j\in\F_q^{N_0}$, with $R_0=1$, and imposes
\begin{equation}
  u_j=\Sigma_{R_j}u+v_j,
  \qquad 0\le j<r.                                  \label{eq:column-relation}
\end{equation}
This relation selects a public subset of possible signatures.  It makes no
assumption about the honest signer and does not change the public verifier.

Before substitution, the verifier feature vector has
$N_{\rm raw}=m_1\ell^2r^2$ labeled coordinates.  Let $z_{\rm ord}(u)$ list
the remaining scalar quadratic forms indexed by ordered pairs of powers of
$S$.  Let $E$ be the public outer map after the fixed permutation that
reconciles the official flattening order with the block-transposed order used
here.  For a fixed relation $R$, let $R_R$ expand $z_{\rm ord}(u)$ into the
raw feature vector, and let $F_{R,v}$ collect the coefficients of the terms
that are linear in $u$.  The substituted verifier is
\begin{equation}
  E\bigl(R_R z_{\rm ord}(u)+F_{R,v}u\bigr)+c_{R,v}=t.
  \label{eq:substituted-verifier}
\end{equation}
Put
\[
 E_R=ER_R,
 \qquad
 \Lambda_{R,v}=EF_{R,v}.
\]
Products of two variable terms give $z_{\rm ord}(u)$.  Products of one
variable term and one offset give the linear term, and products of two offsets
give the constant.

\subsection{How symmetry removes quadratic forms}
For one base form in this restricted scalar map, the ordered label space is
$A\otimes_{\F_q}A$.  The following theorem identifies the pairs that always
give the same polynomial.

\begin{theorem}[Symmetric-square factorization]\label{thm:symmetry-quotient}
Let $A=\F_q[S]$ have dimension $d$, and assume $S^\transpose=S$ and
$P_i^\transpose=P_i$.  For each $i$, the map
\[
  A\otimes_{\F_q}A\longrightarrow \F_q[u]^{\rm hom}_2,
  \qquad
  \xi\otimes\eta\longmapsto u^\transpose\Sigma_\xi P_i\Sigma_\eta u
\]
is invariant under $\xi\otimes\eta\mapsto\eta\otimes\xi$.  It factors
through
\[
  \Sym^2_{\F_q}(A)=
  (A\otimes_{\F_q}A)/\Span\{\xi\otimes\eta-\eta\otimes\xi\}.
\]
Consequently the direct sum over $m_1$ base forms has rank at most
$m_1\binom{d+1}{2}$.  After any public linear map to $M$ outputs, its rank is
at most
\begin{equation}
  \min\left\{M,\;m_1\binom{d+1}{2}\right\}.          \label{eq:symmetry-cap}
\end{equation}
If the minimal polynomial of $S$ has degree $\ell$, then $d=\ell$ and the
ordered feature vector is obtained by duplicating the off-diagonal
coordinates of an unordered feature vector.
\end{theorem}

\begin{proof}
Every element of $A$ is symmetric.  Since the displayed value is a scalar,
\[
 u^\transpose\Sigma_\xi P_i\Sigma_\eta u
 =\bigl(u^\transpose\Sigma_\xi P_i\Sigma_\eta u\bigr)^\transpose
 =u^\transpose\Sigma_\eta P_i\Sigma_\xi u.
\]
Thus every alternating label $\xi\otimes\eta-\eta\otimes\xi$ lies in the
kernel.  The symmetric square has dimension $\binom{d+1}{2}$, and a later
linear map cannot enlarge its image.  When powers of $S$ form a basis, the
coordinate statement is exactly the duplication of each off-diagonal
unordered pair into its two ordered positions.
\end{proof}

The theorem is an upper bound on the common-column scalar image: extra public
dependencies can only reduce that image further.  Every concrete route checks
the exact public rank before solving.  The outer map cannot repair the loss
because, after the relation, it sees only the already-equal polynomial values.
Before the common-column restriction, swapping $(\xi,\eta)$ only transposes
the matrix $U^\transpose\Sigma_\xi P_i\Sigma_\eta U$; it does not make the
matrix entries equal.  The theorem therefore concerns the restricted scalar
verifier, not the unrestricted matrix-valued map.
The factorization itself uses only transposition symmetry.  Odd
characteristic is needed later for the familiar decomposition into symmetric
and alternating tensors and for the polar-form analysis.  The submitted
nonsymmetrized $q=16$ construction is not covered.

\subsection{Mapping the remaining forms to verifier outputs}
Let $z_{\rm sym}(u)$ contain one coordinate for each unordered pair
$0\le a\le b<\ell$ in every base form, and let $D$ duplicate its off-diagonal
coordinates.  By \Cref{thm:symmetry-quotient},
$z_{\rm ord}(u)=Dz_{\rm sym}(u)$.  Define
\begin{equation}
  Q_R=E_RD\in\F_q^{M\times K}.                     \label{eq:QR}
\end{equation}
The columns of $Q_R$ span every output direction reachable by the quotient
quadratics.

\subsection{Separating the quadratic and affine equations}
Choose $W_R$ whose rows form a basis of the left kernel of $Q_R$, and put
\begin{equation}
  C_{R,v}=W_R\Lambda_{R,v}.                          \label{eq:CRv}
\end{equation}
Multiplication by $W_R$ removes every quadratic term and exposes the affine
constraints.  A complementary map $G_R$ keeps the quadratic coordinates.
The two maps together are an invertible change of verifier outputs.

\begin{lemma}[Invertible output split]\label{lem:canonical-residual}
Assume $\rank Q_R=K$.  Choose a public left inverse
$G_R\in\F_q^{K\times M}$ satisfying $G_RQ_R=I_K$.  Then
\[
 \begin{pmatrix}G_R\\W_R\end{pmatrix}\in\F_q^{M\times M}
\]
is invertible.  Writing $b=t-c_{R,v}$, the complete substituted verifier is
equivalent to
\begin{equation}
 C_{R,v}u=W_Rb,
 \qquad
 g_{R,v}(u)=G_Rb,
 \qquad
 g_{R,v}(u)=z_{\rm sym}(u)+G_R\Lambda_{R,v}u.        \label{eq:canonical-split}
\end{equation}
The polar map of $g_{R,v}$ is exactly that of $z_{\rm sym}$.  If $b$ is
uniform in $\F_q^M$, the two right-hand sides in
\eqref{eq:canonical-split} are independent and uniform.
\end{lemma}

\begin{proof}
If a vector $y\in\F_q^M$ is killed by both $G_R$ and $W_R$, then
$W_Ry=0$ implies $y=Q_Ra$ for some $a\in\F_q^K$.  Applying $G_R$ gives
$a=G_RQ_Ra=0$, so $y=0$.  The stacked square matrix is therefore invertible.
Applying its block rows to
$Q_Rz_{\rm sym}(u)+\Lambda_{R,v}u=b$ gives the two displayed systems, and
invertibility gives the converse.  Affine terms have zero polar.  Finally, an
invertible linear transformation sends a uniform vector to a uniform vector
on the product of the two coordinate spaces.
\end{proof}

\begin{proposition}[Exact affine reduction]\label{prop:exact-affine}
Let $\rho_Q=\rank Q_R$, $\lambda=\rank C_{R,v}$, and
$\delta=M-\rho_Q-\lambda$.  Then
\[
  \rank[Q_R\ \Lambda_{R,v}]=\rho_Q+\lambda.
\]
After row reduction, exactly $\delta$ independent affine consistency
conditions remain on the target.  For every target satisfying those
conditions, Gaussian elimination leaves $\rho_Q$ equations of degree at most
two in $N_0-\lambda$ variables, with exactly the same solution set as the
complete substituted verifier.

If $\rho_Q=K$ and $\lambda=M-K$, every target passes the affine consistency
test and the residual system consists of $K$ equations on an affine translate
of $\ker C_{R,v}$.
\end{proposition}

\begin{proof}
The row space of $W_R$ is $\ker Q_R^\transpose$, so
$\ker W_R=(\ker Q_R^\transpose)^\perp=\im Q_R$.  Hence $W_R$ realizes the
quotient map $\F_q^M\to\F_q^M/\im Q_R$.  The image of the columns of
$\Lambda_{R,v}$ in this quotient has dimension
$\rank(W_R\Lambda_{R,v})=\lambda$, proving the first identity.

The projected affine system is $C_{R,v}u=W_Rb$.  Its codomain has dimension
$M-\rho_Q$ and its image has dimension $\lambda$, leaving
$\delta=(M-\rho_Q)-\lambda$ target consistency equations.  Whenever they
hold, parameterize the affine solution space by $N_0-\lambda$ free variables
and substitute into a complementary set of $\rho_Q$ rows.  Every operation is
invertible row reduction or exact affine elimination, so the solution set is
unchanged.  In the full-rank case, $C_{R,v}$ is surjective and $\delta=0$.
\end{proof}

\subsection{Restricting the remaining equations to a square system}
The preceding proposition preserves the full solution set inside the
common-column relation.  The attack now chooses a smaller affine space on
which the number of variables equals the number of remaining equations.  This
second restriction preserves every verifier equation, but it need not contain
a solution.  Root existence is handled separately in
\Cref{sec:sq-full-domain-spectrum}.

\begin{corollary}[Exact square restriction]
\label{cor:exact-square-restriction}
Assume
\[
  \rank Q_R=K,
  \qquad
  \rank C_{R,v}=M-K,
  \qquad
  N_0\ge M.
\]
Then, for every target $t$, choose one solution $u_\star$ to
$C_{R,v}u=W_R(t-c_{R,v})$, a $K$-dimensional subspace
$L\subseteq\ker C_{R,v}$, and an isomorphism
$T:\F_q^K\to L$.  On the affine space $u_\star+L$, the complete substituted
verifier is equivalent to the square system
\begin{equation}
  g_{R,v}(u_\star+Tx)=G_R(t-c_{R,v}),
  \qquad x\in\F_q^K.                               \label{eq:exact-square-restriction}
\end{equation}
Every root reconstructs, through \eqref{eq:column-relation}, a candidate that
satisfies the original verifier equations.  Conversely, every verifier
preimage that satisfies both public restrictions gives a root of
\eqref{eq:exact-square-restriction}.  The corollary does not assert that a root
exists or that it passes the signature-format rejection rule.
\end{corollary}

\begin{proof}
The rank assumption on $C_{R,v}$ makes it surjective, so $u_\star$ exists for
every target.  Its kernel has dimension
\[
  N_0-(M-K)=N_0-M+K\ge K,
\]
so it contains a $K$-dimensional subspace $L$.  The affine equations in
\eqref{eq:canonical-split} hold throughout $u_\star+L$.  The invertible output
change in \Cref{lem:canonical-residual} therefore leaves exactly the $K$
quadratic equations in \eqref{eq:exact-square-restriction}.  That lemma also
gives both directions of the claimed equivalence.
\end{proof}

All nine published parameter shapes satisfy $N_0\ge M$ by
\Cref{tab:parameters}.  Thus, for these shapes, the two displayed rank checks
are enough to construct the exact square system.  They are not enough to show
that this system contains an acceptable root.

\paragraph{Operational meaning.}
The attacker computes every matrix in this section from the supplied public
key.  The attacker may test several public relations and offsets before
solving, but may not replace an unfavorable key.  This paper does not prove
how often that search succeeds or measure its cost separately.  For the
homotopy route those costs are part of the unresolved $\kappa_{\rm hom}$.
For the Just-Guess route, Section~\ref{sec:jg-finite-gate} separately
conditions on the supplied certified $(15,4,7)$ staircase and the fixed-key
rank-$24$ certificate of \Cref{prop:jg-rank24}.  Their deterministic
verification is priced below $2^{50}$ AXN gates, and no fresh-key pass
probability is claimed.  If no tested relation passes the two rank checks, the
conditional recovery results do not apply to that key.  Once the checks pass,
the steps in this section are exact public computations and no verifier
equation is dropped.

\subsection{A Level-I dimension walkthrough}
For $(v,o,\ell,r)=(28,5,4,4)$, we have $n=33$, $N_0=132$, $M=80$,
$m_1=5$, and $K=50$.  The common-column relation leaves $80$ ordered
quadratic labels, but the symmetry quotient leaves only $50$ distinct
quadratics.  In the full-rank affine case, the other $30$ output directions
become linear equations and the full affine residual domain has dimension
$132-30=102$.  The attacker can therefore choose a $50$-dimensional affine
space and obtain $50$ quadratic equations in $50$ variables.  This is the
exact reduction.  Whether the chosen space contains an acceptable root, and
how much work is needed to find one, are the subjects of the conditional
analysis later in the paper.

\section{Preparing square systems for recovery}
\label{sec:stable-rejection}
\label{sec:sq-preparing}

The exact reduction permits any suitable affine slice inside
$\ker C_{R,v}$.  Recovery needs a slice chosen without reading its fresh
coefficients, signatures that pass the format rule, and uniform targets.  This
section provides those inputs; root existence is addressed in
\Cref{sec:sq-full-domain-spectrum}.

\subsection{An unconditional stable kernel}
Generic offsets can supply the missing affine output directions, but their
linear terms may destroy the block grading used by multihomogeneous solvers.
Choose one vector $w\in\F_q^{N_0}$ and multipliers
$\Gamma_0,\ldots,\Gamma_{r-1}\in A$, and set
\begin{equation}
  v_j=\Sigma_{\Gamma_j}w.                            \label{eq:structured-offsets}
\end{equation}
Define the offset row space
\begin{equation}
 U_w=\Span_{\F_q}\{w^\transpose\Sigma_\xi P_i\Sigma_\eta:
     1\le i\le m_1,\ \xi,\eta\in A\}.              \label{eq:Uw}
\end{equation}
For a basis $1,S,\ldots,S^{d-1}$ of $A$, form the complete right-$A$ orbit
\begin{equation}
 O^F_{R,v}=
 \begin{pmatrix}
  F_{R,v}\\F_{R,v}\Sigma_S\\\vdots\\F_{R,v}\Sigma_{S^{d-1}}
 \end{pmatrix},
 \qquad \rho_F=\rank O^F_{R,v}.                    \label{eq:orbit-matrix}
\end{equation}

\begin{theorem}[Stable-lift kernel]\label{thm:stable-kernel}
Assume $S^\transpose=S$, $P_i^\transpose=P_i$, and $A\cong\F_{q^d}$.  Under
\eqref{eq:structured-offsets}:
\begin{enumerate}[label=(\roman*)]
\item $U_w$ is closed under right multiplication by $A$ and
      $\dim_{\F_q}U_w\le m_1d^2$;
\item every row of $F_{R,v}$ and $O^F_{R,v}$ lies in $U_w$;
\item the row space of $O^F_{R,v}$ is an $A$-subspace, so $d\mid\rho_F$;
\item
\[
 H^F_{R,v}=\ker O^F_{R,v}
\]
is $A$-linear, satisfies $F_{R,v}H^F_{R,v}=0$, and has
\begin{equation}
  \dim_A H^F_{R,v}=n-\rho_F/d\ge n-m_1d.             \label{eq:stable-dim}
\end{equation}
\end{enumerate}
No generic orbit-rank equality is required.  An unexpected dependency only
enlarges the stable space.
\end{theorem}

Right $A$-stability puts the orbit rows in an $A$-linear row space, and
rank--nullity gives \eqref{eq:stable-dim}.  The complete proof is in
\Cref{app:slice-proofs}.

Since $C_{R,v}=W_REF_{R,v}$, the stable kernel lies inside
$\ker C_{R,v}$.  Translation along a subspace $L\subseteq H^F_{R,v}$ changes
the quotient quadratics but leaves every offset-linear feature constant.

\subsection{Reserved-line decoupling and the structural preflight}
Let $e_1,\ldots,e_n$ be the native $A$-coordinate basis in which the
key-dependent symmetric public matrices are expanded.  Fix a vinegar index
$j_0$ and use the coordinate-aligned decomposition
\begin{align}
 W&=Ae_{j_0},&
 V'&=\bigoplus_{\substack{j\ {\rm vinegar}\\j\ne j_0}}Ae_j,&
 O_{\rm pub}&=\bigoplus_{j\ {\rm oil}}Ae_j,\notag\\
 V_{\rm pub}&=W\oplus V'\oplus O_{\rm pub},&
 \dim_AW&=1,&
 \dim_AV'&=v-1,\qquad \dim_AO_{\rm pub}=o.
 \label{eq:public-decomposition}
\end{align}
The actual SHAKE stream is deterministic once the public seed is fixed.  For
probability statements only, the \emph{native-coordinate idealized
XOF-transcript model} replaces the relevant bytes by independent uniform bytes
and reduces the native upper-triangular entries of the key-dependent $P_i$
modulo $19$.  This is not a claim about a fixed SHAKE-generated key.  The model
conditions on $E$, $R$, and the full vector
$\Gamma=(\Gamma_0,\ldots,\Gamma_{r-1})$, and invokes independence before any
$A$-linear basis change.  A reduced byte has maximum atom
\[
 \rho_{\rm byte}=\frac{14}{256},
\]
because nine residues have fourteen preimages and ten have thirteen.

\begin{lemma}[Coordinate-aligned reserved-line decoupling]
\label{lem:reserved-decoupling}
Fix $E,R,\Gamma$, choose $0\ne w\in W$, and use the structured offsets
\eqref{eq:structured-offsets}.  Then $Q_R$ depends only on $E$ and $R$, while
$C_{R,v}$, $O^F_{R,v}$, $H^F_{R,v}$, and every deterministically tie-broken
slice selected from them depend only on native entries having at least one
index in $W$.  In the native-coordinate idealized XOF-transcript model they are
independent of the native $V'\times V'$ entries.
\end{lemma}

Every offset contains one coordinate from $W$, so the construction never reads
a native $V'\times V'$ entry.  The complete proof is in
\Cref{app:slice-proofs}.

\begin{lemma}[Maximum-atom pivot bound]\label{lem:max-atom-pivot}
Let $X_1,\ldots,X_N$ be independent finite-field variables, each with maximum
atom at most $\rho_{\rm byte}$.  If a matrix $B$ has rank $t$, then for every
$c$,
\[
 \Pr[BX=c]\le\rho_{\rm byte}^{\,t}.
\]
\end{lemma}

Conditioning on the nonpivot variables leaves one allowed value for each pivot,
so the probabilities multiply.  The complete proof is in
\Cref{app:slice-proofs}.

Put $\mathcal Q_R=\F_q^M/\im Q_R$.  When $\rank Q_R=K$,
$\dim_{\F_q}\mathcal Q_R=M-K$, and $C_{R,v}$ is a coordinate matrix for the
induced linear map $u\mapsto W_R\Lambda_{R,v}u$ to $\mathcal Q_R$.  Let
$\mathbf X_{W,V'}$ be the vector of native independent public-matrix entries
in the $W\times V'$ blocks.  For
$0\ne\lambda\in\mathcal Q_R^\vee$, define $B_\lambda$ by
\[
 \operatorname{vec}\!\left(
   (\lambda\circ C_{R,v})|_{V'}\right)
 =B_\lambda\mathbf X_{W,V'}+b_\lambda.
\]
Here $B_\lambda$ and $b_\lambda$ are fixed once $E,R,\Gamma,\lambda$ and the
native layout are fixed; they do not depend on the sampled entries in
$\mathbf X_{W,V'}$.
Scaling $\lambda$ scales $B_\lambda$, so its rank depends only on
$[\lambda]$.  Put
\[
 t_{\rm src}
 =\min_{0\ne\lambda\in\mathcal Q_R^\vee}\rank B_\lambda.
\]
For an $A$-line $\mathfrak l\subseteq W\oplus O_{\rm pub}$, choose
$0\ne y_{\mathfrak l}\in\mathfrak l$ and write
\[
 O^F_{R,v}y_{\mathfrak l}
 =D_{\mathfrak l}\mathbf X_{\rm nat}+e_{\mathfrak l}
\]
in the complete vector $\mathbf X_{\rm nat}$ of native independent
public-matrix entries.  The coefficient data
$D_{\mathfrak l},e_{\mathfrak l}$ are deterministic before those entries are
sampled, and the rank of $D_{\mathfrak l}$ is independent of the chosen
generator.  Define
\[
 t_W=\rank D_W,\qquad
 t_{\rm oth}=
 \min_{\substack{\mathfrak l\subseteq W\oplus O_{\rm pub}\\
                   \dim_A\mathfrak l=1,\ \mathfrak l\ne W}}
 \rank D_{\mathfrak l}.
\]
These are exact symbolic coefficient-map ranks determined by the fixed
relation, the complete offset pattern, and the native coordinate layout; they
do not depend on the sampled values of the key coefficients.  Define
\begin{align}
 \delta_{\rm src}
 &=\frac{q^{M-K}-1}{q-1}
   \rho_{\rm byte}^{d(v-1)},                          \label{eq:delta-src}\\
 \delta_{\rm proj}
 &=\rho_{\rm byte}^{m_1\binom{d+1}{2}}+
   \left(\frac{|A|^{1+o}-1}{|A|-1}-1\right)
   \rho_{\rm byte}^{m_1d^2},
 \qquad |A|=q^d.                                      \label{eq:delta-proj}
\end{align}

\begin{theorem}[Conditional coefficient-preflight density]
\label{thm:structural-preflight}
Fix the deterministic outer map $E$, a relation $R$, and the complete vector
$\Gamma$.  Assume the exact deterministic checks
\[
\rank Q_R=K,\qquad
t_{\rm src}\ge d(v-1),\qquad
t_W\ge m_1\binom{d+1}{2},\qquad
t_{\rm oth}\ge m_1d^2
\]
pass.  In the native-coordinate idealized XOF-transcript model,
\[
 \Pr[\rank C_{R,v}<M-K]\le\delta_{\rm src},
\]
and
\[
 \Pr[H^F_{R,v}\cap(W\oplus O_{\rm pub})\ne0]
 \le\delta_{\rm proj}.
\]
Consequently full affine rank and injectivity of
$H^F_{R,v}\to V'$ hold jointly with probability at least
$1-\delta_{\rm src}-\delta_{\rm proj}$.  For the six official $\ell=4$
dimension tuples, direct substitution gives
$\delta_{\rm src}+\delta_{\rm proj}<2^{-209}$.  The coefficient-map ranks are
fixed-layout symbolic checks; the realized quotient, affine, and intersection
ranks are exact public checks on a fixed key.
\end{theorem}

The two bounds apply \Cref{lem:max-atom-pivot} to projective output
directions and to $A$-lines in $W\oplus O_{\rm pub}$, then take union bounds.
The complete proof is in \Cref{app:slice-proofs}.

The concrete route records all $r$ entries of $\Gamma$ and runs the displayed
checks.  This release does not assert that every fixed key passes them; the six
$\ell=4$ ledgers are conditional on a successful structural preflight.

\subsection{The guaranteed fresh ordinary square}
Let
\[
  D_0=V'\cap\ker C_{R,v}.
\]
The intersection inequality gives
\begin{equation}
  \dim_{\F_{19}}D_0
  \ge 4(v-1)-(M-K).                                  \label{eq:fresh-slice-dim}
\end{equation}
For all six $\ell=4$ rows this lower bound is at least $K$:
\begin{center}
\small
\begin{tabular}{@{}l r r@{}}
\toprule
Parameters & $\dim D_0$ lower bound & $K$\\
\midrule
$(28,5,4,4)$, $(28,4,4,5)$ & $78$ & $50$\\
$(40,7,4,4)$ & $114$ & $70$\\
$(38,5,4,5)$ & $118$ & $70$\\
$(50,9,4,4)$ & $142$ & $90$\\
$(52,6,4,6)$ & $150$ & $90$\\
\bottomrule
\end{tabular}
\end{center}
Consequently the attacker can choose, from cross data alone, an ordinary
$K$-dimensional slice inside $D_0$.  Its internal $V'\times V'$ coefficients
remain fresh for the spectrum theorem.  This simple dimension fact is the
reason the complete-square route applies to every $\ell=4$ row, including
$(38,5,4,5)$; it does not require the smaller structured stable kernel to have
$A$-dimension $K/4$.

\subsection{Passing the block-symmetry rejection rule}
Version~2.3 rejects a square $\ell=4$ signature if any block is symmetric and
rejects an $\ell=2$ signature if more than $\lfloor n/4\rfloor$ blocks are
symmetric~\cite{SNOVA23}.  Rectangular $\ell=4$ blocks are not subject to the
square-block rule.

For a square $\ell=4$ relation $R=(R_0,R_1,R_2,R_3)$, define the local skew
map
\[
 \mathcal S_R:\F_q^4\longrightarrow\bigwedge^2\F_q^4,
 \qquad
 \mathcal S_R(x)_{a,c}=(R_cx)_a-(R_ax)_c.
\]
For any offsets, one signature block is symmetric exactly when
$\mathcal S_R(x)$ equals a fixed offset-dependent vector.  The public relation
$R_j=S^j$ for the official matrix has $\rank\mathcal S_R=4$.  Thus each
block has at most one rejected common-column value, and the accepted density
in the full common-column space is at least
\begin{equation}
  \alpha_{4,n}=(1-19^{-4})^n.                        \label{eq:l4-acceptance}
\end{equation}
The accepted-root theorem later uses only this global density; it assumes no
independence from the MQ system.

For $\ell=2$, the zero-offset relation used in the main theorem gives a block
\[
 \begin{pmatrix}x&0\\y&0\end{pmatrix},
\]
which is symmetric exactly when $y=0$.  The exact accepted density is derived
in \eqref{eq:sq-l2-alpha}.  A deterministic fixed-skew alternative is recorded
in \Cref{app:slice-proofs}; it is not used by the main theorem.

\subsection{Exactly uniform official targets}
The probability theorems require a uniform target over $\F_{19}^M$, but the
official binary-to-field conversion has a small modulo bias.  Model the target
hash/XOF chunks as independent uniform bit strings, separately from the key
expansion above.  The pinned conversion uses each full eight-byte chunk for
fifteen base-$19$ digits and a final partial chunk for the remainder.  If a
uniform $b$-bit chunk $X$ supplies $a$ digits, retain the message--salt input
only when
\begin{equation}
  X<\left\lfloor\frac{2^b}{19^a}\right\rfloor19^a.   \label{eq:target-rejection-sampling}
\end{equation}
Conditioned on this event, $X\bmod19^a$ is uniform.  Applying the test to each
chunk makes the retained official target exactly uniform, and an invertible
output change preserves uniformity.  Exact evaluation over the nine public
Version~2.3 lengths gives a minimum acceptance probability
$0.1524622985\ldots$, and hence at least $0.152462$.
When target consistency requires more than the $128$-bit salt space, the
existential chosen-message attacker varies a message prefix as well as the
salt.  Target generation occurs before nonlinear solving and adds to, rather
than multiplies, the solve cost.


\section{Root existence on a square slice}
\label{sec:sq-full-domain-spectrum}

The preceding sections produce a square MQ system but do not show that it has
a root with the required format and a nonsingular Jacobian.  This section
derives that probability from an exact collision identity and bounds the
required derivative ranks for the slices used later.

\subsection{Polar forms, rank spectra, and collisions}
For a scalar quadratic polynomial $f$, define the polar form
\[
  \beta_f(x,y)=f(x+y)-f(x)-f(y)+f(0).
\]
In odd characteristic this is bilinear.  For a vector-valued map $g$, a
nonzero $b\in\F_q^K$ selects the scalar quadratic $b\cdot g$.  Its polar rank
measures how many independent directions the quadratic part couples.

For a fixed slice direction $y$, the vector-valued derivative
\begin{equation}
  T_y(x)=g(x+y)-g(x)-g(y)+g(0)
  \label{eq:sq-Ty}
\end{equation}
is linear in $x$.  The value $q^{-\rank T_y}$ is exactly the fraction of its
codomain-dual directions that annihilate $T_y$, after the appropriate
normalization.  Summing these quantities over nonzero $y$ controls the second
factorial moment of the number of roots in a random target--coset pair.  The
same spectrum controls the density of points where the restricted Jacobian
loses rank.

\subsection{The exact first and second moments}

\begin{theorem}[Root-collision identity]\label{thm:collision-identity}
Let $U$ be an $\F_q$-space, let $C:U\to\F_q^c$ be surjective, let
$L\subseteq\ker C$ have dimension $h$, and let
$g:U\to\F_q^K$ have degree at most two.  Sample independent uniform
$d\in\F_q^c$ and $z\in\F_q^K$, then choose a uniform affine coset $\mathcal A$
of $L$ inside $C^{-1}(d)$.  Put
\[
 Z=|\{x\in\mathcal A:g(x)=z\}|,
 \qquad \mu=q^{h-K}.
\]
For $0\ne y\in L$, define
\[
 T_y(x)=g(x+y)-g(x)-g(y)+g(0)
\]
and let $\kappa(y)$ be one if
$-(g(y)-g(0))\in\im T_y$ and zero otherwise.  Then
\begin{align}
 \mathbb EZ&=\mu,                                      \label{eq:moment-first}\\
 \mathbb E[Z(Z-1)]
 &=\mu\sum_{0\ne y\in L}\kappa(y)q^{-\rank T_y}
 \le \mu\sum_{0\ne y\in L}q^{-\rank T_y}.          \label{eq:moment-second}
\end{align}
The domain of $T_y$ is the full space $U$ because the experiment averages over
the affine right-hand side as well as the coset.
\end{theorem}

\begin{proof}
Sampling $d$ and then a uniform $L$-coset in $C^{-1}(d)$ is exactly uniform
over all cosets of $L$ in $U$.  Every point of $U$ therefore lies in the
sampled coset with probability $q^{h-\dim U}$, and its image matches the
independent target with probability $q^{-K}$.  Summing over all points gives
\eqref{eq:moment-first}.

For an ordered pair of distinct roots in one coset, write the second point as
$x+y$ with $0\ne y\in L$.  The collision condition is
\[
 g(x+y)=g(x)
 \quad\Longleftrightarrow\quad
 T_y(x)=-(g(y)-g(0)).
\]
It has $\kappa(y)q^{\dim U-\rank T_y}$ solutions in $x$.  Divide the total
number of ordered collisions by the number
$q^{\dim U-h}q^K$ of target--coset pairs to obtain
\eqref{eq:moment-second}.
\end{proof}

\subsection{Accepted nonsingular roots}

Retain $U,C,g,L$ and the trial from \Cref{thm:collision-identity}, and write
$g=q_2+q_1+q_0$ by homogeneous degree.  For the full-domain derivative $T_y$
from \eqref{eq:sq-Ty}, define
\begin{equation}
 \bar\eta_L^U=\sum_{0\ne y\in L}q^{-\rank T_y}. \label{eq:sq-etaU}
\end{equation}
The superscript records that the trial averages over the affine right-hand side
and the $L$-coset.  The derivative therefore has domain $U$, not $\ker C$.
Restricting it to $\ker C$ gives a different, generally larger quantity used
for fixed-fiber uniqueness, not for this experiment.

Let $X\subseteq U$ be the assignments accepted by the original verifier and
put $\alpha=|X|/|U|$.  A trial samples independent uniform
$d\in\F_q^c$ and $z\in\F_q^K$, and then a uniform affine coset of $L$ inside
$C^{-1}(d)$.

\begin{theorem}[Accepted nonsingular root, full-domain form]
\label{thm:sq-accepted-root}
Assume $\bar\eta_L^U\le\eta$ and
$\alpha>\eta/(q-1)$.  A trial contains an accepted root $x$ satisfying
\[
 \rank D_Lg(x)=h,
 \qquad
 D_Lg(x)[y]=T_y(x)+q_1(y),
\]
with probability at least
\begin{equation}
 p_{\rm ns}
 =q^{h-K}\frac{a^2}{a+\eta},
 \qquad a=\alpha-\frac{\eta}{q-1}.
 \label{eq:sq-pns}
\end{equation}
No independence among acceptance, singularity, the target, and the quadratic
map is assumed.
\end{theorem}

\begin{proof}
Let $Y$ count accepted roots in the sampled target--coset pair at which
$D_Lg$ has rank $h$, and let $Z$ count all roots.  For a projective direction
$[y]\in\mathbb P(L)(\F_q)$, the condition $D_Lg(x)[y]=0$ is an affine system
in $x$ whose linear part is $T_y$.  It therefore holds on at most a
$q^{-\rank T_y}$ fraction of $U$.  Union-bound over projective directions to
obtain
\[
 \frac{|\{x:\rank D_Lg(x)<h\}|}{|U|}
 \le \frac1{q-1}\sum_{0\ne y\in L}q^{-\rank T_y}
 \le \frac\eta{q-1}.
\]
Every point of $U$ is selected by the random fiber--coset with probability
$q^{h-\dim U}$ and matches the independent target with probability $q^{-K}$.
Put $\mu=q^{h-K}$ and $a=\alpha-\eta/(q-1)$.  Hence
$\mathbb EY\ge\mu a$.
The ordered-pair collision identity gives
$\mathbb E[Z(Z-1)]\le\mu\eta$, and
$Y(Y-1)\le Z(Z-1)$.  Therefore
\[
 \mathbb E[Y^2]
 =\mathbb EY+\mathbb E[Y(Y-1)]
 \le \mathbb EY+\mu\eta.
\]
The function $x\mapsto x^2/(x+\mu\eta)$ is increasing for $x\ge0$.
Cauchy--Schwarz therefore gives
\[
 \Pr[Y>0]\ge
 \frac{(\mathbb EY)^2}{\mathbb EY+\mu\eta}
 \ge \mu\frac{a^2}{a+\eta},
\]
which proves \eqref{eq:sq-pns}.  The denominator is strictly smaller than
$1+\eta$ whenever $a<1$.
\end{proof}

The following lemma bounds this spectrum in the idealized transcript model.
Let $\rho_{\rm byte}=14/256$ be the maximum atom of one public-expansion byte
reduced modulo $19$.

\begin{lemma}[Full-domain fresh-spectrum theorem]
\label{lem:sq-full-domain-spectrum}
Assume $q$ is odd and $A=\F_{q^d}$.  Work in the native-coordinate idealized
XOF-transcript model of \Cref{sec:stable-rejection}: the scalar entries in the
fresh symmetric-matrix region are independent reductions of independent uniform
bytes, each having maximum atom $\rho_{\rm byte}=14/256$, while all disjoint
coefficient data are conditioned upon.  Let $L$ be fixed from coefficient data
disjoint from that fresh region.  Suppose every $0\ne y\in L$
has nonzero projection $\bar y$ to a fresh $A$-space $D$, and every nonzero
residual-output combination contains a nonzero coefficient
$c\in\operatorname{Sym}_{\F_q}^2(A)$ in at least one fresh public base form.
If
\[
 \dim_{\F_q}D=d+t,
\]
then for all $0\ne y\in L$ and $0\ne b\in\F_q^K$,
\begin{equation}
 \Pr[b\in\ker T_y^{\mathsf T}]\le\rho_{\rm byte}^t.             \label{eq:sq-pointwise}
\end{equation}
Consequently
\begin{align}
 \mathbb E[\bar\eta_L^U]
 &\le (q^h-1)q^{-K}
 +(q^h-1)(1-q^{-K})\rho_{\rm byte}^t,                            \label{eq:sq-eta-expect}\\
 \Pr\!\left[
  \bar\eta_L^U>(q^h-1)q^{-K}+\varepsilon
 \right]
 &\le
 \frac{(q^h-1)(1-q^{-K})\rho_{\rm byte}^t}{\varepsilon}.
                                                                    \label{eq:sq-eta-tail}
\end{align}
For the reserved-line construction one may take
$D=V'\cong A^{v-1}$ and $t=d(v-2)$.  For the zero-offset construction on the
full public-vinegar summand one may take $D\cong A^v$ and $t=d(v-1)$.
\end{lemma}

For fixed $y\ne0$ and $b\ne0$, the fresh coefficient map has rank at least
$t$.  The maximum-atom bound gives \eqref{eq:sq-pointwise}; expectation and
Markov's inequality give \eqref{eq:sq-eta-expect} and
\eqref{eq:sq-eta-tail}.  The complete proof is in
\Cref{app:slice-proofs}.

For the rest of the paper, use the threshold
\begin{equation}
 \boxed{
 \eta_{128}(h,K)=(q^h-1)q^{-K}+2^{-128}.}
 \label{eq:sq-eta128}
\end{equation}
This allows $2^{-128}$ above the exact full-rank baseline; it does not assume
perfect rank.

\begin{corollary}[Numerical spectrum tails]
\label{cor:sq-spectrum-tails}
For the cost-driving slices, \eqref{eq:sq-eta-tail} with
$\varepsilon=2^{-128}$ gives the following base-two failure exponents.  The
entry in each cell is the weaker official shape when two shapes share a
level.
\begin{center}
\small
\begin{tabular}{@{}c r r r@{}}
\toprule
Route & Level I & Level III & Level V\\
\midrule
$\ell=4$ fast channel core & $-138.12$ & $-237.86$ & $-371.14$\\
$\ell=4$ direct complete square & $-95.64$ & $-178.39$ & $-294.67$\\
$\ell=2$ one-eigenblock core & $-130.17$ & $-263.46$ & $-396.74$\\
$\ell=2$ complete two-block core & $-62.21$ & $-161.50$ & $-260.80$\\
\bottomrule
\end{tabular}
\end{center}
Thus, in that idealized transcript model, even the weakest headline event
fails with probability below $2^{-62}$.
A fixed key can instead publish its exact projective rank histogram, which
certifies \eqref{eq:sq-eta128} without the idealized XOF-transcript model.
\end{corollary}

\section{Conditional recovery for all published parameter sets}
\label{sec:sq-overview}

The exact reduction produces square systems for every published parameter
shape.  Recovery also needs roots with the right format and a solver that
finds them.  We state those conditions once, then give the resulting costs.
The complete-system arithmetic and exact fallback rules are given in
\Cref{app:recovery-details}.

\paragraph{Assumptions.}
For each invoked route, the fixed public key and chosen common-column relation
must pass the exact quotient, affine-rank, and applicable structural checks.
For the $\ell=2$ route, the attacker also retains only targets that pass the
exact public consistency test.  Each invoked slice must satisfy
$\bar\eta_L^U\le\eta_{128}(h,K)$, either under the idealized transcript model
or through an exact spectrum certificate for the fixed key
(\Cref{sec:sq-full-domain-spectrum}).  An exact public determinant test
selects a smaller system when available; if it fails, the algorithm uses the
complete square, so passing this test is not an extra assumption on the key.
The table averages the work of the smaller system and its fallback.  This
average also uses the idealized XOF-transcript bound on how often the
structural and determinant tests fail.  For a fixed key, the exact test
chooses one branch, but the averaged exponent need not describe that key.
Finally, $\mathsf H_{\rm hom}$ must hold for every square system that is
solved.  The Safey--El Din--Schost theorem used here has soft-$O$ complexity;
it does not provide a universal explicit constant for all operations hidden
in our $\kappa_{\rm hom}$.  Accordingly, the displayed work is the arithmetic
ledger multiplied by an unresolved factor $\kappa_{\rm hom}$.

\begin{theorem}[Conditional recovery for all published parameter shapes]
\label{thm:sq-all-nine}
Consider any published $q=19$ Version~2.3 parameter shape under the assumptions
above.  A randomized algorithm solves the certified smaller system when its
exact determinant test passes and otherwise solves the complete square.  It
checks every candidate against all residual equations, confirms that its
coordinates lie in the base field, reconstructs the signature, applies the
signature rejection rule, and runs the original verifier.  It therefore
never returns an invalid signature.

The table gives the arithmetic exponent $e$ before $\kappa_{\rm hom}$ and
the sufficient break-even condition on that unresolved multiplier for the row
to fall below its classical NIST reference.  The Just-Guess analysis is
instantiated only for the Level-I $(48,16,2,2)$ parameters in
Section~\ref{sec:jg-finite-gate}; it is a separate conditional sensitivity,
not a replacement theorem for the other $\ell=2$ rows.
\begin{center}
\small
\begin{tabularx}{\textwidth}{@{}c l >{\raggedright\arraybackslash}X r r@{}}
\toprule
Level & parameters & certified route & exponent $e$ & sufficient break-even\\
\midrule
I   & $(28,5,4,4)$  & smaller channel system, else full square & $128.25$ & $\log_2\kappa_{\rm hom}<14.75$\\
I   & $(48,16,2,2)$ & one eigenblock, else full square & $132.19$ & $\log_2\kappa_{\rm hom}<10.81$\\
I   & $(28,4,4,5)$  & smaller channel system, else full square & $128.25$ & $\log_2\kappa_{\rm hom}<14.75$\\
III & $(40,7,4,4)$  & smaller channel system, else full square & $171.62$ & $\log_2\kappa_{\rm hom}<35.38$\\
III & $(72,24,2,2)$ & complete two-block square & $183.80$ & $\log_2\kappa_{\rm hom}<23.20$\\
III & $(38,5,4,5)$  & smaller channel system, else full square & $171.62$ & $\log_2\kappa_{\rm hom}<35.38$\\
V   & $(50,9,4,4)$  & smaller channel system, else full square & $214.56$ & $\log_2\kappa_{\rm hom}<57.44$\\
V   & $(96,32,2,2)$ & complete two-block square & $233.84$ & $\log_2\kappa_{\rm hom}<38.16$\\
V   & $(52,6,4,6)$  & smaller channel system, else full square & $214.56$ & $\log_2\kappa_{\rm hom}<57.44$\\
\bottomrule
\end{tabularx}
\end{center}
The worst homotopy ledger at each level, before the unresolved multiplier, is
\begin{equation}
 \boxed{132.19+\log_2\kappa_{\rm hom},\qquad
        183.80+\log_2\kappa_{\rm hom},\qquad
        233.84+\log_2\kappa_{\rm hom}}.  \label{eq:sq-headline}
\end{equation}
For the Level-I $(48,16,2,2)$ shape only, the independent Just-Guess analysis
of Section~\ref{sec:jg-finite-gate} gives expected trial cost below
$2^{131.447}$ for the certified diagonal transcript, with its supplied valid $(15,4,7)$ MOR staircase and
fixed-key rank-$24$ certificate.  Its aggregate count follows
all
quadratic branches and all solutions of singular residual linear systems, so
May--Ostuzzi--Ressler Assumption~1 is not used.  The conditional full-rank
joint-pencil event of \Cref{thm:jg-joint-fullrank} fails with probability below
$2^{-114.684}$; on every good Hermitian completion the expected
success-adjusted cost is below $2^{132.447}$ (margin $10.553$ bits below
$143$).  Deterministic verification of the supplied staircase/rank certificate is priced below $2^{50}$ AXN gates; we do not establish its frequency over fresh keys.  These are
Boolean-gate expectation estimates, not wall-clock timings.  Unlike the homotopy rows, however, the Just-Guess route has an explicit depth-first low-space implementation as stated in \Cref{prop:jg-lowspace}.
\end{theorem}

\begin{proof}
The accepted-root theorem of \Cref{thm:sq-accepted-root} gives the root
probability under the stated spectrum condition.  The complete $\ell=4$ and
$\ell=2$ systems are proved in \Cref{sec:sq-complete-square,sec:sq-l2}.
The exact determinant tests, smaller systems, and complete-square fallback are
proved in \Cref{sec:sq-adaptive,sec:sq-l4}.  The final checks rule out invalid
outputs.  Substituting the nine published parameter shapes into
those bounds gives the table; the exact arithmetic rows appear in
\Cref{app:dimensions}.
\end{proof}

\section{Implementation, evidence, and reproducibility}
\label{sec:implementation}

The companion artifact provides formula checkers, one official-key reduction
test, and one toy end-to-end forgery test.  \Cref{tab:evidence-hierarchy}
states what each establishes.  None executes a production-size nonlinear
solver.


Both correspondence tests use a KAT-anchored Python transcription.  The
official-key test compares the substituted map with a separate literal
evaluator but stops before the nonlinear solve.  The toy forger reconstructs
and serializes a signature and checks the rejection rule, but its forger and
verifier share scheme helpers and do not call the official C verifier.  It is
therefore composition evidence rather than an independent implementation.

The release neither tests $\mathsf H_{\rm hom}$ nor measures full wall-clock cost.  Peak memory remains unresolved for the homotopy rows; the separate Level-I Just-Guess route does not share that issue because its enumeration is depth first and its online mutable state is bounded explicitly in \Cref{prop:jg-lowspace}.  The factor $\kappa_{\rm hom}$ is bounded by
an unresolved soft-$O$ and implementation
multiplier; the $\AXN$ ledger
charges the stated multiplication schedule.  \Cref{app:conditional-solver}
records the solver coordinates and homotopy adaptation.
\Cref{sec:jg-finite-gate} supplies a separate Boolean route for the
$\ell=2$ Level-I shape.  For the certified diagonal transcript, using its supplied valid $(15,4,7)$ MOR staircase and fixed-key rank-$24$ certificate, its aggregate
target-averaged
ledger gives expected success-adjusted cost below $2^{132.447}$ on every
Hermitian completion satisfying the full-rank joint-pencil predicate; in the
stated conditional Hermitian-transcript model that predicate fails with
probability below $2^{-114.684}$.  May--Ostuzzi--Ressler Assumption~1 is not
used.  The release includes deterministic verification of the staircase/rank certificate with the conservative $2^{50}$ preflight cap, but does not establish its fresh-key frequency, while
\Cref{sec:low-state-xl} does not prove its production solving degree.

A production fixed-key record would need the exact preflights, a certified
slice spectrum, the complete solver and filtering transcript, the reconstructed
signature, and the original verifier transcript.  Only an instrumented run can measure production wall-clock time and cache/memory traffic.  This caveat should not be read as an exponential-space requirement for the Level-I Just-Guess route: \Cref{prop:jg-lowspace} gives a direct low-space implementation, while the rank-$24$ preflight itself has the conservative $64$ MB mutable-state bound in Section~\ref{sec:jg-finite-gate}.

\begin{table}[t]
\centering
\caption{Evidence shipped with the release.  ``Included'' means that the
named script and its pinned input are present, not that the headline recovery
cost has been measured.}
\label{tab:evidence-hierarchy}
\small
\renewcommand{\arraystretch}{1.14}
\begin{tabularx}{0.98\textwidth}{@{}>{\raggedright\arraybackslash}p{0.19\textwidth}>{\raggedright\arraybackslash}p{0.16\textwidth}X@{}}
\toprule
Evidence & Availability & What it establishes\\
\midrule
Formula and ledger checks & Included & Numerical consequences of the dimensions, probability models, and cost conventions supplied to the scripts; not scheme correspondence or solver execution.\\
Scalar circuit netlist & Included & The emitted $150$-gate $\F_{19}$ multiplication netlist is correct on all $361$ valid inputs.  The extension-field figures are recurrence counts, not emitted netlists.\\
Official KAT verification and reduction & Included for one Level-I $(28,5,4,4)$ key & A Python transcription parses the compressed public key, accepts the supplied signature, rejects an altered message, and obtains quotient/affine ranks $50/30$.  The resulting $50$-equation system in $102$ variables is left unsolved.\\
Reduced zero-offset forgery & Included at unofficial $(2,1,2,2)$ & The complete restricted map is interpolated; its image equals the explicit rank-three quotient image, leaving one target-consistency coordinate before public slice search.  The main case, eight committed fresh-key regressions, and a separate 24-forgery/100-rank stress check pass in the shared KAT-anchored Python transcription.  This is toy composition evidence, not independent C-verifier or production-cost evidence.\\
Level-I $\ell=2$ MOR staircase census & Included for all $100$ archived official KAT public keys & The deterministic omit-form-0 $(15,4,7)$ MOR construction succeeds on $100/100$ independently expanded serialized public keys.  This is finite-corpus evidence, not a fresh-key probability theorem.\\
Certified Level-I $(15,4,7)$ MOR/rank-$24$ preflight & Included for one official key & Exact reconstruction reproduces the archived transformed system coefficient-for-coefficient; the supplied staircase is valid and the full sixteen-form polar pencil satisfies the maximum-rank identity $\rank_A K_{\rm full}(\lambda)=24$ for every nonzero $\lambda$, by the exact $680$-check certificate of \Cref{prop:jg-rank24}.\\
Legacy off-diagonal activation check & Included for the same key & All $2016$ native coordinate-pair restrictions are checked exactly; $2015$ have rank~$3$, and exhaustive enumeration of the unique rank-$2$ pair checks all $129{,}600$ non-sign collisions.  This is an independent fixed-key cross-check; the new headline uses the stronger conditional joint-pencil theorem instead.\\
Fixed-schedule Level-I Boolean audit & Included & The artifact regenerates the $<2^{21}$ aggregate ledger, the explicit $4\times4$ RREF/kernel count, the sequential-check cost, the conditional Hermitian block atoms, and the numerical joint-pencil/second-moment bounds used in \Cref{prop:jg-level1-gates}.\\
Production-size solver or forgery & Not included or claimed & No empirical wall-clock measurement of the Level-I Just-Guess attack, no fresh-key frequency for the staircase/rank-$24$ certificate, and no production-size forgery transcript.\\
\bottomrule
\end{tabularx}
\end{table}

\FloatBarrier

\section{Necessary dimension floors within the common-column model}
\label{sec:upgrade-repair}

We give two necessary output floors for redesigns that retain the
common-column architecture.  The first removes a square slice from every
target fiber.  The stronger second floor removes the corresponding
zero-offset chosen-message slice.  Neither floor is sufficient for security.

\begin{proposition}[Dimensionally square common-column slice bound]
\label{prop:upgrade-fresh-square-floor}
Let $A=\F_{q^d}$ and retain the public vinegar summand
$V'$ of base-field dimension $d(v-1)$.  Suppose an exact affine reduction has
$M'$ outputs, $K'$ quotient quadratics, and full affine rank $M'-K'$.  Then
\[
 \dim_{\F_q}(V'\cap\ker C_{R,v})
 \ge d(v-1)-(M'-K').
\]
If $M'\le d(v-1)$, dimensions alone leave an ordinary $K'$-dimensional slice
carrying all $K'$ residual equations.  Equality of equation and variable
counts does not establish a root, zero-dimensionality, nonsingularity, or an
efficient solve.  Eliminating this dimensional guarantee requires
$M'>d(v-1)$.  Under the coupling $M'=dx'$, this forces $x'\ge v$ and
$M'\ge dv$.
\end{proposition}

\begin{proof}
Use
\[
 \dim(V'\cap\ker C_{R,v})
 \ge\dim V'-\rank C_{R,v}
 =d(v-1)-(M'-K').
\]
The right side is at least $K'$ exactly when $M'\le d(v-1)$.  Negate and
round to the next multiple of $d$.
\end{proof}

Let $K_{\rm cc,ord}=\min\{M,m_1d^2\}$ be the pre-symmetrization ordered-label
cap within the same common-column model.  Retaining symmetric forms requires
$m_1'\ge\lceil K_{\rm cc,ord}/\binom{d+1}{2}\rceil$ and
$M'\ge K_{\rm cc,ord}$.  Under
$M'=dx'$ and $m_1'=\lceil x'/d\rceil$, all three necessities combine as
follows.

\begin{theorem}[Combined necessary common-column output floor]
\label{thm:upgrade-combined-floor}
Any redesign that retains symmetric public forms, the public vinegar count,
the public-vinegar affine geometry above, full affine rank $M'-K'$, and the
Version~2.3 output/base-form coupling, while both restoring
$K_{\rm cc,ord}$ and eliminating the dimensionally guaranteed fresh slice,
must satisfy
\begin{equation}
 x'\ge\max\left\{
 v,
  \left\lceil\frac{K_{\rm cc,ord}}d\right\rceil,
 d\left(
  \left\lceil\frac{K_{\rm cc,ord}}{\binom{d+1}{2}}\right\rceil-1
 \right)+1
 \right\},
 \qquad M'=dx'.                                      \label{eq:upgrade-combined-floor}
\end{equation}
Across the nine official rows, this requires $45.00\%$--$60.71\%$ more
outputs (\Cref{thm:upgrade-zero-offset-floor}).
\end{theorem}

\begin{proof}
The second and third terms in \eqref{eq:upgrade-combined-floor} are the two
ordered-cap requirements under $m_1'=\lceil x'/d\rceil$.  The term $v$ is
\Cref{prop:upgrade-fresh-square-floor}.  Their maximum is therefore necessary.
Substitution gives the affine-fiber column in the comparison below.
\end{proof}

\paragraph{Zero-offset targets.}
\label{sec:upgrade-zero-offset-repair}

The preceding floor assumes a full-rank affine lift.  A zero-offset attacker
can instead filter chosen-message targets through the quadratic image and use
the full common-column domain.  Removing this square slice requires a stronger
floor.

\begin{proposition}[Zero-offset dimensionally square restriction]
\label{prop:upgrade-zero-offset-square-persistence}
Consider any symmetric redesign for which the zero-offset common-column
relation with $R_0=1$ is defined.  Let $Q_R$ be its quadratic output matrix
and put $\rho_Q=\rank Q_R$.  Public row reduction gives an invertible output
change under which the zero-offset verifier is equivalent to
\begin{equation}
 g_{\rho_Q}(u)=z\in\F_q^{\rho_Q},
 \qquad
 t_{\rm con}=0\in\F_q^{M'-\rho_Q},                     \label{eq:upgrade-general-zero-split}
\end{equation}
where $g_{\rho_Q}$ consists of $\rho_Q$ independent quadratic output
combinations.  For an exactly uniform target, the consistency event in
\eqref{eq:upgrade-general-zero-split} has probability $q^{-(M'-\rho_Q)}$ and,
conditioned on it, $z$ is uniform in $\F_q^{\rho_Q}$.

With zero offsets no reserved line is needed, so use the full public-vinegar
space $V_{\rm vin}\cong A^v$, of base-field dimension $dv$.  If
\begin{equation}
 \rho_Q\le dv,                                          \label{eq:upgrade-zero-square-condition}
\end{equation}
then dimensions alone leave an ordinary $\rho_Q$-dimensional slice
$L\subseteq V_{\rm vin}$ carrying all $\rho_Q$ residual equations: it has
$\rho_Q$ equations in $\rho_Q$ variables.  This dimension equality alone
proves neither root existence nor efficient recovery.  Consequently, ruling
out this chosen-message dimensional configuration requires
\begin{equation}
 \rho_Q>dv.                                             \label{eq:upgrade-zero-square-rank-floor}
\end{equation}
\end{proposition}

\begin{proof}
With zero offsets, the substituted verifier is $Q_Rz_{\rm sym}(u)=t$ and
contains no affine or constant term.  Choose a row basis of $Q_R$ and extend
it to an invertible basis change on the $M'$ output coordinates.  The first
$\rho_Q$ transformed rows are the independent residual quadratics
$g_{\rho_Q}$; the remaining rows vanish identically on the quadratic image and
therefore impose the consistency equations in
\eqref{eq:upgrade-general-zero-split}.  An invertible linear transformation
sends a uniform target to independent uniform residual and consistency
coordinates, proving the distribution statement.  Under
\eqref{eq:upgrade-zero-square-condition}, choose any $\rho_Q$-dimensional
linear subspace of $V_{\rm vin}$ and restrict all $\rho_Q$ residual equations
to an affine coset of it.  This gives the asserted dimensionally square
restriction.
\end{proof}

Symmetry bounds the actual quotient rank by
\begin{equation}
 \rho_Q\le \min\left\{M',\;m_1'\binom{d+1}{2}\right\}.
                                                               \label{eq:upgrade-zero-square-cap}
\end{equation}
Hence making \eqref{eq:upgrade-zero-square-rank-floor} even possible forces
both terms on the right of \eqref{eq:upgrade-zero-square-cap} above $dv$.

\begin{theorem}[Output floor eliminating the zero-offset dimension guarantee]
\label{thm:upgrade-zero-offset-floor}
Retain symmetry, the public vinegar count, and the Version~2.3 coupling
\[
 M'=dx',\qquad m_1'=\left\lceil\frac{x'}d\right\rceil.
\]
Put
\[
 C_d=\binom{d+1}{2},
 \qquad
 m_{\rm sq}=\left\lfloor\frac{dv}{C_d}\right\rfloor+1.
\]
A necessary condition for eliminating the zero-offset chosen-message
dimension guarantee is
\begin{equation}
 x'\ge
 \max\left\{v+1,\;d(m_{\rm sq}-1)+1\right\},
 \qquad M'=dx'.                                        \label{eq:upgrade-zero-offset-floor}
\end{equation}
If the redesign also restores the common-column ordered-label cap
$K_{\rm cc,ord}$, the
right side of \eqref{eq:upgrade-zero-offset-floor} must additionally be
maximized with
\[
  \left\lceil\frac{K_{\rm cc,ord}}d\right\rceil,
 \qquad
  d\left(\left\lceil\frac{K_{\rm cc,ord}}{C_d}\right\rceil-1\right)+1.
\]
For the nine official rows, the zero-offset condition dominates the ordered-cap
terms.  The two repair floors compare as follows.
\begin{center}
\small
\begin{tabular}{@{}c l r r r@{}}
\toprule
&&& \multicolumn{2}{c}{necessary outputs $M'$}\\
\cmidrule(l){4-5}
Level & parameters & current $M$ & affine fiber & zero offset\\
\midrule
I   & $(28,5,4,4)$  & $80$  & $116\;(45.00\%)$ & $180\;(125.00\%)$\\
I   & $(48,16,2,2)$ & $64$  & $96\;(50.00\%)$ & $130\;(103.125\%)$\\
I   & $(28,4,4,5)$  & $80$  & $116\;(45.00\%)$ & $180\;(125.00\%)$\\
III & $(40,7,4,4)$  & $112$ & $180\;(60.71\%)$ & $260\;(132.14\%)$\\
III & $(72,24,2,2)$ & $96$  & $144\;(50.00\%)$ & $194\;(102.08\%)$\\
III & $(38,5,4,5)$  & $100$ & $152\;(52.00\%)$ & $244\;(144.00\%)$\\
V   & $(50,9,4,4)$  & $144$ & $228\;(58.33\%)$ & $324\;(125.00\%)$\\
V   & $(96,32,2,2)$ & $128$ & $192\;(50.00\%)$ & $258\;(101.5625\%)$\\
V   & $(52,6,4,6)$  & $144$ & $228\;(58.33\%)$ & $324\;(125.00\%)$\\
\bottomrule
\end{tabular}
\end{center}
The zero-offset floor requires $101.5625\%$--$144.00\%$ more outputs.
\end{theorem}

\begin{proof}
Condition \eqref{eq:upgrade-zero-square-rank-floor} and the two caps in
\eqref{eq:upgrade-zero-square-cap} require
$dx'>dv$ and $m_1'C_d>dv$.  The first inequality is equivalent to
$x'\ge v+1$.  The second requires $m_1'\ge m_{\rm sq}$, and the least integer
$x'$ with $\lceil x'/d\rceil\ge m_{\rm sq}$ is
$d(m_{\rm sq}-1)+1$.  This proves
\eqref{eq:upgrade-zero-offset-floor}.  The two extra ordered-cap terms are
the same feature and output necessities used in
\Cref{thm:upgrade-combined-floor}.  Direct substitution gives the table.
\end{proof}

The stronger result is a dimension-only theorem.  A redesign below the table
can retain the square system and rely on the consistency factor
$q^{M'-\rho_Q}$, its solving cost, or a changed target interface.  The floor
applies only if the redesign instead removes the square slice while preserving
the common-column architecture and parameter coupling.

\section{Conclusion}\label{sec:conclusion}
The common-column restriction exposes a structural weakness in the symmetric
$q=19$ SNOVA verifier.  Symmetry forces the restricted ordered-pair features
to factor through a symmetric square, and an invertible public output change
separates the surviving quadratics from the affine and target-consistency
conditions.  For all nine Version~2.3 shapes satisfying the stated public rank
checks, the result is an exact square quadratic system that preserves the
original verifier and rejection rule.

The strongest concrete result is for the Level-I $\ell=2$ shape.  One official
KAT key carries a deterministic $(15,4,7)$ staircase and a maximum rank-$24$
diagonal-pencil certificate.  Together with target-averaged enumeration and
the conditional joint-pencil theorem, this gives expected success-adjusted
work below $2^{132.447}$ Boolean gates on every good Hermitian completion.
The modeled probability of a bad Hermitian completion, conditioned on the
diagonal transcript, is below $2^{-114.684}$.  The attack uses no
May--Ostuzzi--Ressler Assumption~1, its deterministic preflight is below
$2^{50}$ gates, and its online enumeration is explicitly low-space.  The
remaining limitations are equally explicit: the diagonal certificate has no
proved fresh-key frequency, and no production wall-clock forgery is claimed.

For the remaining rows, the structural reduction is still exact but the
numerical recovery estimates retain the homotopy hypothesis and the unresolved
multiplier $\kappa_{\rm hom}$.  The redesign bounds show that merely restoring
the old ordered-label count is insufficient if the common-column architecture
and current parameter coupling are retained.

\clearpage
\bibliographystyle{unsrtnat}
\begingroup
\footnotesize
\setlength{\bibsep}{0pt}
\Urlmuskip=0mu plus 1mu\relax
\sloppy
\bibliography{references}
\endgroup

\clearpage
\appendix
\section{Conditional solver coordinates and finite-field homotopy}
\label{app:conditional-solver}

The exact reduction in Section~4 needs only the dimension of the symmetric
square.  The conditional recovery algorithms use the more structured
coordinates and finite-field homotopy adaptation recorded here.

\subsection{Frobenius coordinates for the symmetric square}
For $A=\F_{q^d}$ and $j>0$, define
\[
 \phi_j([x\otimes y]_{\rm sym})=xy^{q^j}+x^{q^j}y,
 \qquad
 \phi_0([x\otimes y]_{\rm sym})=xy.
\]
If $d$ is even and $j=d/2$, the value lies in $\F_{q^{d/2}}$.

\begin{theorem}[Frobenius-distance normal form]\label{thm:frobenius-normal}
There is an $\F_q$-linear isomorphism
\[
 \Sym^2_{\F_q}(\F_{q^d})\cong
 \begin{cases}
   A^{(d+1)/2},& d\text{ odd},\\
   A^{d/2}\oplus\F_{q^{d/2}},& d\text{ even},
 \end{cases}
\]
given by the maps $\phi_j$ for the distinct Frobenius distances.  In
particular,
\begin{align}
 \Sym^2_{\F_q}(\F_{q^4})&\cong
    \F_{q^4}\oplus\F_{q^4}\oplus\F_{q^2},            \label{eq:l4-channels}\\
 \Sym^2_{\F_q}(\F_{q^2})&\cong
    \F_{q^2}\oplus\F_q.                              \label{eq:l2-channels}
\end{align}
\end{theorem}

\begin{proof}
Extend scalars to an algebraic closure and identify
$A\otimes_{\F_q}\overline{\F}_q$ with a product of $d$ copies indexed by
$\mathbb Z/d\mathbb Z$.  Frobenius cyclically shifts the primitive
idempotents.  The map $\phi_0$ records the $d$ diagonal unordered pairs.
For $0<j<d/2$, the conjugates of $\phi_j$ record the orbit of unordered pairs
$\{i,i+j\}$, an orbit of size $d$.  If $d$ is even and $j=d/2$, the pair is
repeated after a half-turn, leaving an orbit of size $d/2$.  These orbits
partition all unordered pairs.  The displayed map is therefore a coordinate
permutation after scalar extension and hence an isomorphism over $\F_q$.
\end{proof}

For $\ell=4$, the three coordinate groups have highest homogeneous forms of
ordinary degrees $2$, $q+1=20$, and $q^2+1=362$ over $A=\F_{19^4}$.
For $\ell=2$, the quotient becomes two groups with bidegrees $(2,0)$,
$(1,1)$, and $(0,2)$ after scalar extension.

\subsection{Finite-cardinality symbolic homotopy}
The selected-degree-free branch adapts a multihomogeneous homotopy algorithm.
The scope of the adaptation is important.

Partition $N=\sum_{j=1}^b n_j$ variables into $b$ blocks
$X_1,\ldots,X_b$.  Let $f_1,\ldots,f_N$ be nonconstant polynomials represented
by a straight-line program of length $L$, and suppose $f_i$ has
multidegree at most $\mathbf d_i=(d_{i,1},\ldots,d_{i,b})$.  Define
\begin{align}
 B=C_{\mathbf n}(\mathbf d)
 &= [\vartheta_1^{n_1}\cdots\vartheta_b^{n_b}]
    \prod_{i=1}^N\left(\sum_{j=1}^b d_{i,j}\vartheta_j\right),
                                                        \label{eq:bg-Bezout}\\
B^+=C'_{\mathbf n}(\mathbf d')
 &=\operatorname{coeffsum}\!\left(
    \prod_{i=1}^N\left(\vartheta_0+
      \sum_{j=1}^b d_{i,j}\vartheta_j\right)
    \bmod\langle\vartheta_0^2,
      \vartheta_1^{n_1+1},\ldots,\vartheta_b^{n_b+1}\rangle
    \right).                                         \label{eq:bg-Bplus}
\end{align}

\paragraph{Inherited homotopy hypothesis $\mathsf H_{\rm hom}$.}
For the input system and the chosen random parameters, assume that the
incidence, deformation-curve regularity, specialization, and cleaning
hypotheses used by Proposition~5, Algorithm~2, and Remark~15 of
Safey El Din and Schost remain valid after the two finite-field substitutions
made below~\cite{SafeyElDinSchost2018}.  In particular, the relevant homotopy
components must have the expected dimensions and the specializations and
inversions used during Newton lifting must avoid their exceptional loci.  This
paper proves the start-system and separator substitutions below; it does not
derive $\mathsf H_{\rm hom}$ from the specific ideals produced by \SNOVA{}, and
no released script tests it at a published parameter set.

\begin{theorem}[Conditional finite-cardinality homotopy adaptation]
\label{thm:finite-cardinality-homotopy}
Assume $\mathsf H_{\rm hom}$.  Let $k$ be a perfect finite field and put
\[
 E_0=\max\left\{
      \max_j\sum_i d_{i,j},\;8B^2
      \right\}.
\]
If $|k|\ge E_0$, a randomized symbolic-homotopy algorithm whose separating
form has a coefficient vector sampled uniformly from $k^N$ returns a
zero-dimensional parametrization containing every nonsingular root with
probability at least $7/8$.  On every run it returns a subset of the
nonsingular roots or reports failure.  Its arithmetic cost is
\[
 \widetilde O\!\left(
 B B^+
 \left(L+\sum_{i,j}d_{i,j}+N^2\right)N
 \right)
\]
operations in $k$.  This soft-$O$ theorem does not instantiate its implied
constant or polylogarithmic factors.
More generally, if $|k|\ge\max_j\sum_i d_{i,j}$ and $|k|>B^2$, the same
algorithm succeeds with probability at least $1-B^2/|k|$.
\end{theorem}

\begin{proof}
Proposition~5 and Algorithm~2 of Safey El Din and Schost give the underlying
homotopy construction under their stated hypotheses, with the bad-run behavior
described in their Remark~15~\cite{SafeyElDinSchost2018}.
The statement here is a new finite-cardinality adaptation, not a verbatim
corollary of that proposition.  Its two replacements for the cited
characteristic-dependent choices are proved next.

For block $j$, put $D_j=\sum_i d_{i,j}$ and choose distinct
$\alpha_{j,0},\ldots,\alpha_{j,D_j-1}\in k$.  The affine moment-curve forms
\[
 \kappa_{j,u}(X_j)=X_{j,1}+\alpha_{j,u}X_{j,2}+
   \cdots+\alpha_{j,u}^{n_j-1}X_{j,n_j}+\alpha_{j,u}^{n_j}
\]
have the same Vandermonde properties as the start forms in the cited proof:
any $n_j$ are independent and any $n_j+1$ are inconsistent.  The product
start system can be written explicitly.  Put
\[
 s_{i,j}=\sum_{\nu=1}^i d_{\nu,j},\qquad s_{0,j}=0,
\]
and, for $1\le i\le N$, set
\[
 g_i(X)=
 \prod_{j=1}^b\ \prod_{u=s_{i-1,j}}^{s_{i,j}-1}
 \kappa_{j,u}(X_j).
\]
The factor ranges are disjoint.  An isolated common zero chooses, for each
$i$, one vanishing factor of $g_i$; exactly $n_j$ choices must come from
block $j$.  The coefficient in \eqref{eq:bg-Bezout} counts these choices,
including the $d_{i,j}$ possible factors.  The Vandermonde property gives one
point for each choice, while inconsistency of any $n_j+1$ forms prevents
additional vanishing factors.  The corresponding block-diagonal linear
Jacobian is nonsingular, so all $B$ roots are simple.  They are enumerated
within the same soft-linear cost.  The first term in $E_0$ supplies enough
distinct nodes.

The homotopy proof next requires a well-separating linear form.  It identifies
at most $B^2$ nonzero vectors on which the coefficient vector must not vanish.
The cited proof samples from a one-parameter moment-curve family, which is why
its field-size condition includes an additional factor depending on $N$.
Here the coefficient vector is instead sampled uniformly from the full space
$k^N$.  For each fixed nonzero vector, its annihilator inside $k^N$ is a proper
$k$-linear subspace (possibly of codimension greater than one), so a uniform
coefficient vector annihilates it with probability at most $|k|^{-1}$.  Since
$|k|\ge8B^2$, a union bound gives success at least $7/8$.

Under $\mathsf H_{\rm hom}$, the remaining Newton lifting, specialization, squarefree cleaning, and
zero-dimensional parametrization steps of Algorithm~2 are algebraic over a
perfect field and use only field operations and inversion at nonsingular
points.  Finite fields are perfect, so those correctness and arithmetic
arguments carry over.  Direct substitution into the original system and a
Jacobian-rank test remove every spurious or singular output; hence a run
returns only nonsingular roots or reports failure.  As in the cited
Remark~15, a bad run may return a strict subset, and completeness is not
certified at comparable cost.  The attack wrapper therefore treats absence of
an accepted root as a restart rather than as a proof that none exists.
\end{proof}

The theorem bounds time, not memory.  A parametrization can contain $O(B)$
field elements.  The soft-$O$ factors, fixed transforms, factorization,
filtering, memory traffic, and implementation constants are included later in
the route-specific multiplier $\kappa_{\rm hom}$.

\section{Detailed recovery instantiations}
\label{app:recovery-details}

This appendix derives the arithmetic ledger behind
\Cref{thm:sq-all-nine}.  It first fixes one cost model, then derives the
complete systems used as fallbacks and the certified smaller systems used by
the optimized route.  The final subsection supplies the coefficient-level
proofs deferred from \Cref{sec:sq-preparing,sec:sq-full-domain-spectrum}.

\subsection{Arithmetic cost model}
\label{sec:sq-arithmetic}

All routes below use the same ledger.  Pointwise extension-field products use
the tower basis and the displayed multiplication recurrences.  The factor
$\kappa_{\rm hom}$ contains basis conversion, bulk evaluation and
interpolation transforms, additions, inversions, factorization, control,
memory, and all other hidden costs.  Thus every displayed exponent is a
multiplication-schedule estimate, not a wall-clock or low-memory bound.  The
artifact contains a narrower variable-product count, but no theorem uses it,
and it performs no production-size solve.

\paragraph{Status of $\kappa_{\rm hom}$.}
The finite-cardinality homotopy theorem used above inherits a soft-$O$
complexity statement from Safey--El Din--Schost.  That notation suppresses
polylogarithmic factors and constants; the cited result does not provide a
universal Boolean-gate factor that can be inserted into the present ledgers.
Moreover, our $\kappa_{\rm hom}$ also includes basis conversion, inversions,
factorization, filtering, control, indexing, communication, and memory costs.
Expanding only multipoint evaluation or Newton lifting therefore cannot prove
a bound on the full multiplier.  We consequently make no numerical claim on
$\kappa_{\rm hom}$.  The tables instead report break-even thresholds that a
future complete implementation or explicit circuit analysis would have to
establish.

An explicit finite-arithmetic expansion is possible in principle: the
homotopy is univariate in the deformation parameter $t$, Schost's formal
Newton operator is precision-doubling, and the dominant step is ordinary
polynomial arithmetic at degrees $B$ and $B^+$.  Instantiating a concrete
multiplication algorithm (e.g.\ Cantor--Kaltofen specialized to characteristic
19) and summing the Newton stages would replace the soft-$O$ by a numerical
factor.  For the critical Level-I $\ell=2$ row the available head-room is
only $10.81$ bits; for Levels III and V the head-rooms are $23.2$ and
$38.2$ bits, so those rows are far more tolerant of residual polylog factors.
We do not carry out that expansion here.  For the Level-I $\ell=2$ shape the
paper instead uses the independent Just-Guess Boolean ledger of
Section~\ref{sec:jg-finite-gate}.  That ledger removes the MOR random-tree
assumption by aggregate target averaging; its remaining qualifications are the
supplied certified $(15,4,7)$ staircase and fixed-key rank-$24$ certificate, plus
the stated conditional Hermitian-transcript model.

\begin{theorem}[Constructive scalar netlist and tower recurrences]
\label{thm:sq-circuits}
The circuit model assigns unit cost to two-input AND, XOR, and XNOR gates,
with free constants, wires, and fan-out.  For canonical five-bit
representatives, the following circuit-size assignments are constructive:
\[
\begin{array}{c|rrrr}
 &\text{addition}&\text{negation}&\text{subtraction}&\text{multiplication}\\
\hline
\F_{19}   &42&16&58&150\\
\F_{19^2} &84&32&116&692\\
\F_{19^4} &168&64&232&2628.
\end{array}
\]
The emitted $150$-gate scalar netlist is correct on all $19^2=361$ canonical
input pairs.  The $692$ and $2628$ extension-field multiplication figures are
constructive composition counts from the displayed primitive recurrences;
the release does not emit complete netlists for those two multipliers.  The
tower representations are
\[
 \F_{19^2}=\F_{19}[u]/(u^2+1),\qquad
 \F_{19^4}=\F_{19^2}[v]/(v^2-(1+u)).
\]
They are exactly isomorphic to the manuscript basis
$\F_{19}[t]/(t^4-t-1)$ via
\[
 u=1+2t+15t^2+5t^3,\qquad
 v=4+13t+13t^2+t^3,
\]
and the basis $(1,u,v,uv)$ has determinant $16\ne0$ modulo $19$.
\end{theorem}

\begin{proof}
The scalar circuit first converts a canonical input
$x=(x_0,\ldots,x_4)$ into a sign bit $s$ and a four-bit magnitude $m$.  On
the valid inputs $0,\ldots,18$, set
\[
 a=x_1x_3,\quad b=x_2x_3,\quad
 s=a\oplus b\oplus ax_2\oplus x_4,
\]
\[
 (m_0,m_1,m_2,m_3)
 =(x_0\oplus s,x_1\oplus s,x_2,x_3\oplus b).
\]
A truth-table check gives $s=0,m=x$ for $x\le9$ and $s=1,m=19-x$ otherwise.
The two decoders use $18$ gates.  The magnitude product and output sign reach
gate $80$.  The product is at most $81$, so write it as $\ell+32h$ with
$h\in\{0,1,2\}$, replace $32h$ by $13h\bmod19$, and reduce once.  This reaches
gate $130$.  The final $20$ gates map the result to $r$ or $-r\bmod19$ as
required by the sign.  The released evaluator checks all $361$ canonical
input pairs against the canonical output.

For $\F_{19^2}$, Karatsuba uses three scalar products and two input sums.  The
shared output negations give
\[
 5\cdot42+2\cdot16+3\cdot150=692.
\]
For $\F_{19^4}$, three $\F_{19^2}$ products cost $3\cdot692$.  The two input
sums cost $2\cdot84$.  Multiplication by $1+u$ costs one scalar addition and
one subtraction.  The two outputs account for the remaining extension-field
operations, giving
\[
 2\cdot84+3\cdot692+(42+58)+84+(84+116)=2628.
\]
Direct polynomial arithmetic verifies both tower equations, the displayed
coordinates, and the basis determinant.
\end{proof}

\begin{lemma}[Conservative well-separating-form success]
\label{lem:sq-separator}
For a square multihomogeneous core with B\'ezout number $B$, represented over a
finite field $E$ large enough for the finite-cardinality start system, the
well-separating-form construction has a set of at most $B^2$ nonzero forbidden
coefficient vectors.  A uniform linear form therefore satisfies every
separation condition with probability at least
\[
 \sigma_E(B)=1-\frac{B^2}{|E|}.
\]
Conditional on this event and $\mathsf H_{\rm hom}$, the finite-cardinality
homotopy algorithm returns a zero-dimensional parametrization containing every
nonsingular core root, with one-sided error.
\end{lemma}

\begin{proof}
Failure means that the coefficient vector of the linear form annihilates one
of at most $B^2$ nonzero vectors.  A uniform coefficient vector annihilates
each fixed vector with probability $|E|^{-1}$, so the union bound gives
$B^2/|E|$.  Pairwise separation alone would count fewer hyperplanes, but it
does not cover every condition required by the homotopy parametrization.
\end{proof}

For a core over $A\in\{\F_{19^2},\F_{19^4}\}$ with B\'ezout number $B$ and
start-node requirement $D_{\rm start}$, let $g_A$ be $692$ or $2628$ and define
\begin{equation}
 \Gamma_A(B,D_{\rm start})=
 \min_{\substack{e\ge1:\ |A|^e\ge D_{\rm start},\
                   |A|^e>B^2\\2e-1\le|A|}}
 \frac{(2e-1)g_A}{1-B^2/|A|^e}.                      \label{eq:sq-gamma}
\end{equation}
Evaluation/interpolation at $2e-1$ base-field points gives the number of
charged pointwise multiplications in the numerator; the fixed transforms
themselves remain in $\kappa_{\rm hom}$.  \Cref{lem:sq-separator} gives the
denominator.  The numerical artifact
exhausts all admissible $e$.


\subsection{Complete-system fallback}
\label{sec:sq-complete-route}

The fallback solves every residual equation.  The $\ell=4$ shapes use an
ordinary square over a fresh public slice, while the $\ell=2$ shapes use all
three bidegree families.  Both routes apply the accepted-root theorem to the
Jacobian of the complete system and filter every returned point through the
original verifier.

\subsubsection{The ordinary square for \texorpdfstring{$\ell=4$}{ell=4}}
\label{sec:sq-complete-square}

\begin{theorem}[Conditional recovery from the complete square system]
\label{thm:sq-direct-square}
Retain the setting of \Cref{thm:sq-accepted-root}, assume
$\mathsf H_{\rm hom}$, and suppose $h=K$.
Restrict the complete residual map $g$ to the public affine $K$-dimensional
slice and solve all $K$ equations together.  On the event
$\bar\eta_L^U\le\eta$, a trial contains an accepted root at which the
Jacobian of this very system is nonsingular with probability at least
\[
 p_\square=\frac{a^2}{a+\eta},
 \qquad a=\alpha-\frac{\eta}{q-1}.
\]
If every equation has ordinary degree at most two, then
\[
 B_\square=2^K,
 \qquad
 B_\square^+=2^K\left(1+\frac K2\right).
\]
Let
\[
 D_K=\binom{K+1}{2},\qquad
 L_\square=D_K+2K(D_K+K+1),\qquad
 H_\square=L_\square+2K+K^2.
\]
Under $\mathsf H_{\rm hom}$ and over any field satisfying the displayed
finite-cardinality conditions, the multiplication work before
$\kappa_{\rm hom}$ is
\begin{equation}
 W_\square=
 \frac{\Gamma_A(2^K,2K)}{p_\square}
 2^{2K}\left(1+\frac K2\right)H_\square K.
 \label{eq:sq-square-work}
\end{equation}
The algorithm is Las Vegas after complete public filtering and assumes no
selected Macaulay degree, exact corank, unique root, preferred core, or orbit
sweep.
\end{theorem}

\begin{proof}
Because $h=K$, \Cref{thm:sq-accepted-root} applies to the Jacobian of the
complete residual system itself.  No equation projection is needed.  The
ordinary B\'ezout number of $K$ quadratics in $K$ variables is $2^K$.  For the
companion coefficient, taking no auxiliary term contributes $2^K$, while
taking it from one of the $K$ factors contributes $2^{K-1}$.  Therefore
$B_\square^+=2^K+K2^{K-1}$.

The evaluation program forms the $D_K$ quadratic monomials and then evaluates
all $K$ affine quadratics.  The conditional homotopy bound and
\Cref{lem:sq-separator} now give \eqref{eq:sq-square-work}.  Complete public
filtering checks every returned point against the full residual system,
affine reconstruction, signature rule, and original verifier.
\end{proof}

\begin{corollary}[Conditional direct-square routes for the six $\ell=4$ shapes]
\label{cor:sq-l4-direct-square}
For every official $\ell=4$ shape, a public key that passes the exact quotient
and affine-rank checks leaves a fresh ordinary $K$-dimensional slice in the
public-vinegar region.  Assume the applicable spectrum inequality and
$\mathsf H_{\rm hom}$.  Applying \Cref{thm:sq-direct-square}, embedding the
$\F_{19}$ coefficients into $\F_{19^2}$, and charging the constructive
$692$-gate multiplication recurrence gives normalized exponents
\[
 \boxed{142.60696,\qquad185.45088,\qquad227.56041}
\]
on the worse shape at Levels I, III, and V.  The corresponding nominal
headroom below $143,207,272$ is
\[
 \boxed{0.39304,\qquad21.54912,\qquad44.43959}
\]
bits before $\kappa_{\rm hom}$.  The two direct Level-I comparisons survive
only for multipliers below approximately $1.313157$ and $1.313691$,
respectively.  A factor of two erases both crossings.
\end{corollary}

\begin{proof}
The fresh-slice inequality gives
$\dim(V'\cap\ker C_{R,v})\ge K$ for all six rows, so the required ordinary
$K$-slice exists.  The full-domain spectrum theorem applies with $t=4(v-2)$
and threshold $\eta_{128}(K,K)$.  Substitute $K=50,70,90$ in
\eqref{eq:sq-square-work}.  Exhaustive search over the separator extension
degree selects $e=13,17,22$ over $\F_{19^2}$.  The ledger then conditions on
the exact structural check, divides by the certified spectrum probability,
and uses the conservative signature-acceptance bound.
\end{proof}

\subsubsection{The zero-offset split for \texorpdfstring{$\ell=2$}{ell=2}}
\label{sec:sq-l2}

For the three official $\ell=r=2$ rows put
\[
 A=\F_{19^2},\qquad m=m_1=o,\qquad M=4m,\qquad K=3m.
\]
Use the zero-offset column relation
\begin{equation}
 u_0=u,\qquad u_1=0.                                  \label{eq:sq-l2-zero}
\end{equation}

\begin{proposition}[Exact zero-offset target split]
\label{prop:sq-zero-split}
Let $Q_R$ be the quotient-output matrix for \eqref{eq:sq-l2-zero}.  The
route applies to a fixed key only after the exact public test
$\rank Q_R=K$.  No probability claim replaces that fixed-key test for the
three $\ell=2$ rows.  After a passing test, choose a left inverse $G_R$
and a full-row-rank matrix
$W_R$ spanning the left kernel of $Q_R$.  The complete substituted verifier
is equivalent to
\begin{equation}
 W_Rt=0,\qquad z_{\rm sym}(u)=G_Rt.                  \label{eq:sq-l2-target-split}
\end{equation}
For a uniform official target, consistency has probability
$19^{-(M-K)}=19^{-m}$.  Conditioned on consistency, the residual target is
uniform in $\F_{19}^K$, so the attacker can filter message and salt inputs
before solving.  Charging $2^{32}$ AXN operations per retained uniform target
gives the additive exponents
\[
 16\log_2 19+32=99.96684,
 \quad24\log_2 19+32=133.95026,
 \quad32\log_2 19+32=167.93368.
\]
\end{proposition}

\begin{proof}
Zero offsets remove the affine and constant terms, leaving
$Q_Rz_{\rm sym}(u)=t$.  The invertible output change with block rows
$G_R,W_R$ gives \eqref{eq:sq-l2-target-split}.  It preserves uniformity and
separates the residual target from the consistency coordinates.  Since target
filtering precedes the nonlinear solve, its work adds to the solve ledger.
\end{proof}

In one block, \eqref{eq:sq-l2-zero} has matrix
$\left(\begin{smallmatrix}x&0\\y&0\end{smallmatrix}\right)$ and is symmetric
exactly when $y=0$.  The $n=v+o$ block conditions use disjoint coordinates.
Thus the exact density accepted by the Version~2.3 rule is
\begin{equation}
 \alpha_n=19^{-n}\sum_{j=0}^{\lfloor n/4\rfloor}
          \binom nj18^{n-j}.                          \label{eq:sq-l2-alpha}
\end{equation}
Choose all solver slices inside the full public-vinegar summand $A^v$,
independently of its fresh internal symmetric coefficients.  By
\Cref{lem:sq-full-domain-spectrum}, the fresh exponent is
$t=2(v-1)$.

\paragraph{The complete two-block system.}

Take an $A$-linear slice of dimension $s=3m/2$.  Its base-field dimension is
$h=2s=K$, and the complete quotient system has $m$ equations in each
bidegree family
\[
 (2,0),\qquad(1,1),\qquad(0,2).
\]
It is therefore square and includes every residual equation.

\begin{theorem}[Conditional complete two-block $\ell=2$ recovery]
\label{thm:sq-l2-complete}
On the event
$\bar\eta_L^U\le\eta_{128}(K,K)$, a consistent target--coset trial contains
an accepted full-Jacobian root with probability at least
\begin{equation}
 p_{\rm comp}^{(2)}
 =\frac{a_2^2}{a_2+\eta_{128}(K,K)},
 \qquad
 a_2=\alpha_n-\frac{\eta_{128}(K,K)}{18}.            \label{eq:sq-l2-root}
\end{equation}
The exact multihomogeneous quantities are
\begin{align}
 B_m^{(2)}&=2^{2m}\binom m{m/2},                       \label{eq:sq-l2-B}\\
 (B_m^{(2)})^+
 &=B_m^{(2)}\left(1+2m+\frac{m^2}{m+2}\right).       \label{eq:sq-l2-Bplus}
\end{align}
Assuming $\mathsf H_{\rm hom}$, finite-cardinality homotopy over
$\F_{19^2}$ recovers the roots without choosing a Macaulay degree or a
preferred subsystem.  Frobenius descent and complete public filtering reject
every invalid candidate.  Including the additive target filter and the
$2^{-128}$ spectrum normalization gives the following exponents before
$\log_2\kappa_{\rm hom}$:
\[
 132.97067,\qquad183.79463,\qquad233.84399
\]
for Levels I, III, and V.
\end{theorem}

\begin{proof}
Apply \Cref{thm:sq-accepted-root} with $h=K$ and the accepted density from
\eqref{eq:sq-l2-alpha}.  Over the splitting field, each base form contributes
one equation to each diagonal family and one to the cross family.  Hence the
B\'ezout number is
\[
 [t_0^st_1^s](2t_0)^m(t_0+t_1)^m(2t_1)^m.
\]
The cross factor must contribute $m/2$ powers to each block, proving
\eqref{eq:sq-l2-B}.

For the companion coefficient, the term with no auxiliary variable contributes
$B_m^{(2)}$.  Deleting a diagonal factor contributes
$2^{2m-1}(\binom m{m/2}+\binom m{m/2+1})$, while deleting a cross factor
contributes $B_m^{(2)}$.  Summing and using
$\binom m{m/2+1}/\binom m{m/2}=m/(m+2)$ proves
\eqref{eq:sq-l2-Bplus}.  Finally, \eqref{eq:sq-gamma} accounts for the
separating form, and every later operation is an exact filter.
\end{proof}

\Needspace{8\baselineskip}
For the complete two-block system, define
\begin{align}
 D_s&=\binom{s+1}{2},& C_s&=(s+1)^2,\notag\\
 L_{\rm comp}&=2D_s+s^2+4m(D_s+s+1)+2mC_s,\notag\\
 H_{\rm comp}&=L_{\rm comp}+6m+K^2.
\end{align}
On a good spectrum event, the multiplication work before
$\kappa_{\rm hom}$ is
\begin{equation}
 W_{\rm comp}^{(2)}
 =\frac{\Gamma_{\F_{19^2}}(B_m^{(2)},2s)}
        {p_{\rm comp}^{(2)}}
   B_m^{(2)}(B_m^{(2)})^+H_{\rm comp}K.              \label{eq:sq-l2-work}
\end{equation}
Divide by the exact spectrum-event density and add the target-generation
ledger.

\subsubsection{All-nine complete-system costs}

\begin{corollary}[All-nine complete-system fallback]
\label{thm:sq-all-nine-simple}
Consider any published $q=19$ Version~2.3 parameter shape.  Suppose the public
key and column relation pass the exact rank checks, the complete-system bound
$\bar\eta_L^U\le\eta_{128}(K,K)$ holds, and $\mathsf H_{\rm hom}$ holds for
the invoked system.  Then the complete-system algorithm checks every candidate
against the affine reconstruction, signature rule, and original verifier, so
it never returns an invalid signature.  The algorithm does not choose a
Macaulay degree or assume a unique root, predicted corank, or predetermined
Jacobian minor.

The worst direct costs at Levels I, III, and V are
\begin{equation}
 \boxed{142.60696,\qquad185.45088,\qquad233.84399}.   \label{eq:sq-simple-headline}
\end{equation}
The nine direct rows and their optimized counterparts appear together in
\Cref{tab:exact-ledger}.  These values include the idealized target-generation
ledger from \Cref{sec:sq-preparing} and precede the unresolved multiplier
$\kappa_{\rm hom}$.
\end{corollary}

\begin{proof}
For $\ell=4$, \Cref{cor:sq-l4-direct-square} gives the ordinary squares.  For
$\ell=2$, \Cref{thm:sq-l2-complete} gives the two-block squares.  Both use
\Cref{thm:sq-accepted-root}, the separating-form bound, and complete public
filtering.  Substitution of the nine published shapes gives
\Cref{tab:exact-ledger}; taking the largest direct entry at each level gives
\eqref{eq:sq-simple-headline}.
\end{proof}

\subsection{Certified smaller systems and optimized costs}
\label{sec:sq-l4}

The complete systems above give the fallback.  A smaller route uses fewer
linear combinations of the residual equations.  The algorithm first checks
the leading Jacobian determinant with public data.  It solves the smaller
system only when this exact test passes and otherwise returns to the complete
system.

\subsubsection{Certified projection and complete-square fallback}
\label{sec:sq-adaptive}

The public test checks that the highest homogeneous part of the selected
system's Jacobian determinant is not the zero polynomial.  In the applications
below, evaluating that polynomial at one fixed public point suffices.

\begin{theorem}[Certified smaller system with complete-square fallback]
\label{thm:sq-adaptive}
Fix a public construction for which a common structural and spectrum event
$\mathcal S$ makes the complete-system fallback valid.  Suppose:
\begin{enumerate}
\item when the exact public determinant test passes, the smaller system has
expected ledger cost at most $W_f$;
\item when the test fails, the complete fallback has expected work at most
$W_F$;
\item under the idealized XOF-transcript distribution, the test fails with
probability at most $\delta$, while
$\Pr[\mathcal S^c]\le\varepsilon$.
\end{enumerate}
For every fixed key in $\mathcal S$, the algorithm tests the smaller system
and otherwise uses the complete fallback.  Under the idealized transcript
distribution, its conditional expected work obeys
\begin{equation}
 \mathbb E[W\mid\mathcal S]
 \le W_f+
 \min\!\left\{1,\frac\delta{1-\varepsilon}\right\}
 (W_F-W_f)_+ .                                      \label{eq:sq-adaptive-work}
\end{equation}
If an independent exact structural test succeeds with probability
$p_{\rm str}$, the transcript-model-normalized ledger divides the right side by
$p_{\rm str}(1-\varepsilon)$.  Without independence one may instead use any
proved lower bound on $\Pr[\mathcal S]$.
\end{theorem}

\begin{proof}
For a fixed key, the determinant test is a certificate rather than a
probabilistic premise.  A passing key uses the smaller system, while a failing
key uses the complete system.  Both branches apply the full public filters.

For the transcript-model cost, let $F$ be the event that the determinant test
fails.  Then
\[
 \Pr[F\mid\mathcal S]
 \le\frac{\Pr[F]}{\Pr[\mathcal S]}
 \le\min\{1,\delta/(1-\varepsilon)\}.
\]
If $q=\Pr[F\mid\mathcal S]$ and $W_F\ge W_f$, the branch average is
$W_f+q(W_F-W_f)$.  If $W_F<W_f$, always using the fallback only lowers the
cost.  This gives the positive-part expression.  The last statement divides
by the density of idealized transcripts that reach the solver.
\end{proof}

\subsubsection{The smaller Level-I \texorpdfstring{$\ell=2$}{ell=2} system}

For the $\ell=2$ rows, retain $A=\F_{19^2}$, $m=o$, $K=3m$, and the accepted
density $\alpha_n$ from \eqref{eq:sq-l2-alpha}.  For $s\le m$, select $s$
diagonal equations and solve one eigenblock of
$A^s$.  A descended full root reconstructs the other block by Frobenius.
The core is a square system of $s$ ordinary quadratics, with
\[
 B_s=2^s,\qquad B_s^+=2^s(1+s/2).
\]
At a fixed nonzero public point, the fresh leading row of each selected
public form is surjective onto $(A^s)^\vee$.  Sequential exposure therefore
gives
\begin{equation}
 p_{\rm Jac}^{(2)}(s)
 \ge\prod_{j=1}^s(1-\rho_{\rm byte}^{2j})>0.9970003329.          \label{eq:sq-l2-jac}
\end{equation}
A nonzero leading determinant has degree at most $s$ over $A$, so its singular
point density is at most $s/19^2$.  Put
$\eta_s=\eta_{128}(2s,3m)$ and $a_s=\alpha_n-s/19^2$.  The singleton and
second-moment arguments then give the following lower bound for an accepted
root at which this core is nonsingular:
\begin{equation}
 p_s=19^{2s-3m}
 \max\left\{
   a_s-\eta_s,
   \frac{a_s^2}{a_s+\eta_s}
 \right\}.                                           \label{eq:sq-l2-diag-root}
\end{equation}

At Level I, set $s=m=16$ and run the exact leading-determinant test.  A passing
key uses the one-eigenblock system.  A failing key uses the complete two-block
system from \Cref{thm:sq-l2-complete}.  The fallback theorem gives normalized
exponent
\begin{equation}
 \boxed{132.18807}.                                   \label{eq:sq-l2-I-adaptive}
\end{equation}
The diagonal-only and complete-only exponents in the reported ledger are
$132.18928$ and $132.97067$.  The adaptive value is slightly smaller than
the first rounded value because the ledger applies the common spectrum
normalization once and combines the two branches before taking logarithms.
At Levels III and V, the complete system is already faster.

The separate Just Guess sensitivity for the Level-I $\ell=2$ row appears in
\Cref{sec:jg-finite-gate}.  It is not used by the main recovery theorem.

\subsubsection{The smaller \texorpdfstring{$\ell=4$}{ell=4} channel systems}

Let $A=\F_{19^4}$ and let $L\cong A^s$ be the public slice.  Under the
Frobenius-distance normal form, each public base form supplies an $A$-valued
degree-two channel and an $A$-valued degree-$20$ channel.  Write their highest
homogeneous parts as
\[
 F_{i,0}^{\rm top}(x)=C_{i,0}(x,x),\qquad
 F_{i,1}^{\rm top}(x)=C_{i,1}(x,x^{19}).
\]

\begin{lemma}[Fresh channel decomposition and row surjectivity]
\label{lem:sq-channel-surj}
Restriction of a fresh symmetric $\F_{19}$-bilinear form to the injected
public-vinegar copy of $L$, followed by the Frobenius-distance transform, is
surjective onto
\[
 \operatorname{Sym}_A^2(L^\vee)
 \oplus\operatorname{Bil}_A(L,L^{(19)};A).
\]
For every fixed $0\ne x_\star\in L$, the map from one public form's fresh
coefficients to the pair of leading Jacobian rows
\[
 \bigl(DF_{i,0}^{\rm top}(x_\star),
       DF_{i,1}^{\rm top}(x_\star)\bigr)
 \in L^\vee\oplus L^\vee
\]
is surjective over $A$.
\end{lemma}

\begin{proof}
After extending scalars, the four diagonal blocks form one Frobenius orbit.
They descend to an arbitrary symmetric $A$-bilinear form.  The four adjacent
blocks form a disjoint orbit and descend to an arbitrary bilinear form between
$L$ and its Frobenius twist.  Since the two coefficient sets are disjoint, the
map onto the two channels is surjective.

Formal differentiation gives
\[
 DF_{i,0}^{\rm top}(x_\star)[h]=2C_{i,0}(x_\star,h),
 \qquad
 DF_{i,1}^{\rm top}(x_\star)[h]=C_{i,1}(h,x_\star^{19}).
\]
Evaluation at nonzero $x_\star$ is surjective for arbitrary bilinear forms.
For the symmetric channel, given $\lambda\in L^\vee$, choose
$\phi(x_\star)=1$ and use the symmetric form
\[
 \frac12(\phi\otimes\lambda+\lambda\otimes\phi
          -\lambda(x_\star)\phi\otimes\phi).
\]
Its evaluation at $x_\star$ is $\lambda$, which proves row surjectivity.
\end{proof}

Choose $a$ whole coordinates of the degree-two channel and $b$ whole
coordinates of the degree-$20$ channel, with $a+b=s$.  The determinant
test is exact on the public key.  The following transcript-model bound does
not assume that an invertible linear change preserves independence under the
biased reduction modulo $19$.

\begin{lemma}[Transcript-model failure bound for the Jacobian test]
\label{lem:sq-l4-projective-jacobian}
At the fixed public test point, let $J\in A^{s\times s}$ be the selected
leading-Jacobian matrix, where $A=\F_{19^4}$.  Then
\begin{equation}
 \Pr[\det J=0]
 \le \frac{|A|^s-1}{|A|-1}\rho_{\rm byte}^{4s}.      \label{eq:sq-l4-jac}
\end{equation}
For the cost-minimizing values $s=10,14,18$, the failure bounds are at most
$3.55350\cdot10^{-5}$, $6.56029\cdot10^{-5}$, and
$1.21113\cdot10^{-4}$, respectively.
\end{lemma}

\begin{proof}
If $J$ is singular, its right kernel contains a projective direction
$[y]\in\mathbb P^{s-1}(A)$.  Fix $[y]$.  Row surjectivity implies that forcing
all $s$ selected rows to annihilate $y$ imposes $s$ independent equations over
$A$.  These become $4s$ independent affine equations over $\F_{19}$ in fresh
public symbols.  The maximum-atom bound gives probability at most
$\rho_{\rm byte}^{4s}$ for this direction.

There are $(|A|^s-1)/(|A|-1)$ projective directions.  A union bound proves
\eqref{eq:sq-l4-jac}.  Although weaker than the full-rank product for uniform
field elements, this bound applies to the biased product distribution in the
idealized transcript model.
\end{proof}

Conditional on passing the determinant test, put
\[
 a_{a,b}=\alpha-\frac{a+19b}{19^4}.
\]
The sharpened core-root probability is
\begin{equation}
 p_{a,b,s}
 =19^{4s-K}\frac{a_{a,b}^2}
                    {a_{a,b}+\eta_{128}(4s,K)}.       \label{eq:sq-l4-root}
\end{equation}
The exact B\'ezout data for this core are
\[
 B=2^a20^b,\qquad B^+=B(1+a/2+b/20).
\]
Define its straight-line evaluation cost by
\begin{align*}
 L_{a,b,s}
 &=\mathbf1_{a>0}\left[\binom{s+1}{2}
   +2a\left(\binom{s+1}{2}+s+1\right)\right]\\
 &\quad+\mathbf1_{b>0}\left[6s+s^2+2b(s^2+2s+1)\right],
\end{align*}
Under $\mathsf H_{\rm hom}$, the multiplication ledger for a key that passes
the determinant test is
\begin{equation}
 \frac{\Gamma_{\F_{19^4}}(B,2a+20b)}{p_{a,b,s}}
 BB^+\bigl(L_{a,b,s}+2a+20b+s^2\bigr)s.             \label{eq:sq-l4-fast-work}
\end{equation}
Exhaustive enumeration of all $35$, $63$, and $99$ admissible profiles gives
\[
 (s;a,b)=(10;5,5),\quad(14;7,7),\quad(18;9,9)
\]
at Levels I, III, V.  Conditional on the exact structural test and the
stated spectrum/model events, and assuming $\mathsf H_{\rm hom}$, the
worse-shape smaller-system exponents before $\log_2\kappa_{\rm hom}$ are
$126.260$, $167.038$, and $207.273$.

The artifact regenerates every numerical input for these profiles.  The paper
makes no numerical claim for a sixteen-orbit or output-size-optimized branch.

\subsubsection{Combining the certified branches}
\label{sec:sq-all-nine-details}

For $\ell=2$, the Level-I test selects one eigenblock when its determinant
certificate passes and otherwise uses the complete two-block square.  The
complete square is already cheaper at Levels III and V.  For $\ell=4$, the
projective Jacobian bound of \Cref{lem:sq-l4-projective-jacobian} controls the
fast branch, and a failed test returns to the complete system of
\Cref{cor:sq-l4-direct-square}.  The artifact combines the two branch costs
before taking logarithms, charges the union of their spectrum events, and
normalizes only by the probability of the common spectrum event.

After expanding the soft-$O$ bound explicitly, we retain the resulting
implementation multiplier $\kappa_{\rm hom}$.  Comparing the arithmetic
ledgers with the classical NIST references gives
the row-specific sufficient thresholds in \Cref{thm:sq-all-nine}; no claim is
made here that those thresholds hold.  Separately, for the Level-I $(48,16,2,2)$ shape, the Boolean Just-Guess
route gives expected success-adjusted cost below $2^{132.447}$ on every
Hermitian completion satisfying the conditions of Section~\ref{sec:jg-finite-gate}: the supplied certified $(15,4,7)$ target-independent MOR staircase and fixed-key rank-$24$ certificate,
and the conditional full-rank joint-pencil event of \Cref{thm:jg-joint-fullrank}.  Its aggregate target-averaged ledger
does not use May--Ostuzzi--Ressler Assumption~1.

\subsection{Proofs for square-slice preparation and spectrum bounds}
\label{app:slice-proofs}

This subsection collects the coefficient-level proofs used in
\Cref{sec:sq-preparing,sec:sq-full-domain-spectrum}.  The corresponding
statements and proof ideas remain in the main body.

\begin{proof}[Proof of \Cref{thm:stable-kernel}]
The family in \eqref{eq:Uw} has $m_1d^2$ members.  It is stable under right
multiplication by $A$ because, for every $\zeta\in A$,
\[
 (w^\transpose\Sigma_\xi P_i\Sigma_\eta)\Sigma_\zeta
 =w^\transpose\Sigma_\xi P_i\Sigma_{\eta\zeta},
\]
which is another member of the span.  Every offset cross term contains one
copy of $w$.  If $w$ occurs on the right, symmetry and commutativity move it to
the left.  Thus every row of $F_{R,v}$ lies in $U_w$, and right stability puts
every orbit row there as well.

The complete right orbit in \eqref{eq:orbit-matrix} makes its row space an
$A$-vector space, so $d\mid\rho_F$.  Its first block is $F_{R,v}$, hence its
kernel kills every offset-linear feature.  If $x$ is in the kernel and
$\zeta\in A$, each $S^a\zeta$ is an $\F_q$-linear combination of the chosen
basis powers.  The same orbit rows therefore kill $\Sigma_\zeta x$, proving
that the kernel is $A$-linear.  Rank-nullity over $A$ gives
\eqref{eq:stable-dim}.
\end{proof}

\begin{proof}[Proof of \Cref{lem:reserved-decoupling}]
Every structured offset lies in $W$, so each offset-linear coefficient
contains one copy of $w$.  Symmetry moves that copy to the left when needed.
The coefficient therefore reads only native matrix entries with at least one
index in $W$.  Right multiplication by $A$ preserves this coordinate block,
as do the deterministic orbit and slice constructions.

The native $V'\times V'$ entries come from disjoint XOF symbols and remain
unread.  They are therefore independent of the selected slice in the stated
model.  This argument must stay in the native coordinates because the biased
byte-reduction distribution is not invariant under an arbitrary $A$-linear
basis change.
\end{proof}

\begin{proof}[Proof of \Cref{lem:max-atom-pivot}]
Choose $t$ pivot columns.  Condition on all nonpivot variables.  Expose the
pivot variables in echelon order; at each step at most one value can satisfy
the next independent equation, and that value has probability at most
$\rho_{\rm byte}$.  Multiplication gives the claim.
\end{proof}

\begin{proof}[Proof of \Cref{thm:structural-preflight}]
First consider failure of the affine source rank.  If
$\rank C_{R,v}<M-K$, some nonzero
$\lambda\in\mathcal Q_R^\vee$ annihilates $C_{R,v}$.  In particular,
$(\lambda\circ C_{R,v})|_{V'}=0$.  For fixed $[\lambda]$, this is an affine
system of rank at least $d(v-1)$ in the independent native variables
$\mathbf X_{W,V'}$; \Cref{lem:max-atom-pivot} bounds its probability by
$\rho_{\rm byte}^{d(v-1)}$.  Scalar multiples define the same event, and
$\mathbb P(\mathcal Q_R^\vee)(\F_q)$ has
$(q^{M-K}-1)/(q-1)$ points.  A union bound proves
\eqref{eq:delta-src}.

Next consider failure of the projection to $V'$.  By
\Cref{thm:stable-kernel}, $H^F_{R,v}$ is $A$-linear.  A nonzero intersection
with $W\oplus O_{\rm pub}$ must therefore contain an $A$-line $\mathfrak l$.
For $\mathfrak l=W$, the coefficient-rank hypothesis and
\Cref{lem:max-atom-pivot} give probability at most
$\rho_{\rm byte}^{m_1\binom{d+1}{2}}$.  For every other $A$-line the
analogous bound is $\rho_{\rm byte}^{m_1d^2}$.  The $(1+o)$-dimensional
$A$-space $W\oplus O_{\rm pub}$ contains
$(|A|^{1+o}-1)/(|A|-1)$ lines, exactly one of which is $W$; a union bound
gives \eqref{eq:delta-proj}.  Finally, the kernel of the projection
$H^F_{R,v}\to V'$ is precisely
$H^F_{R,v}\cap(W\oplus O_{\rm pub})$.

A final union bound combines the source-rank and projection failures without
assuming independence.  Direct substitution gives the stated numerical
bound.
\end{proof}

\begin{proof}[Proof of \Cref{lem:sq-full-domain-spectrum}]
Fix $y\ne0$ and $b\ne0$.  Choose a fresh public base form whose corresponding
symmetric-square coefficient $c$ is nonzero.  Writing
$c=\sum_\nu a_\nu\odot b_\nu$, the fresh symmetric-matrix coefficient in the
scalar derivative $b\cdot T_y(x)$ is
\[
 \Xi_{c,\bar y}(x)=
 \sum_\nu\bigl((a_\nu\bar y)\odot(b_\nu x)
 +(b_\nu\bar y)\odot(a_\nu x)\bigr).
\]
We first bound the kernel of this coefficient map.  If $x\notin A\bar y$,
project onto the cross summand between the disjoint lines $A\bar y$ and $Ax$.
After identifying both lines with $A$, this projection is the ordinary
symmetrization of $c$.  It is nonzero because symmetrization is injective on
$\operatorname{Sym}^2(A)$ in odd characteristic.  Hence the kernel is
contained in $A\bar y$, and the coefficient map has rank at least
$\dim_{\F_q}D-d=t$.

The event $b\in\ker T_y^{\mathsf T}$ forces a rank-$t$ affine system in
independent fresh symbols.  Conditioning on the nonpivot symbols and applying
the maximum-atom bound to the pivots proves \eqref{eq:sq-pointwise}.

Finally, the dual-kernel identity gives
\[
 q^{-\rank T_y}
 =q^{-K}|\ker T_y^{\mathsf T}|
 =q^{-K}\left(1+\sum_{b\ne0}
   \mathbf1_{\{b\in\ker T_y^{\mathsf T}\}}\right).
\]
Its expectation is at most
$q^{-K}+(1-q^{-K})\rho_{\rm byte}^t$.  Summing over $y\ne0$ proves
\eqref{eq:sq-eta-expect}.  The excess above the baseline
$(q^h-1)q^{-K}$ is nonnegative, so Markov's inequality proves
\eqref{eq:sq-eta-tail}.  The two stated choices of $D$ have dimensions
$d(v-1)$ and $dv$.
\end{proof}

\paragraph{An unused fixed-skew format check.}
A deterministic alternative is also available:
for the official matrix
\[
 S=\begin{pmatrix}1&2\\2&15\end{pmatrix},
 \qquad
 R_1=9I+10S=\begin{pmatrix}0&1\\1&7\end{pmatrix},
\]
choose offsets with $(v_1)_{b,0}-(v_0)_{b,1}=1$ in every block.  Then each
block has fixed nonzero skew $1$, so every root passes the rejection rule.
This fixed-skew route is a useful separate algebraic check but is not needed
for the reported recovery theorem.

\section{A Boolean Level-I attack via a certified 15-form Just-Guess redesign}
\label{sec:jg-finite-gate}

The homotopy ledger retains an unresolved multiplier $\kappa_{\rm hom}$.  For
its numerically weakest Level-I row, $(v,o,\ell,r)=(48,16,2,2)$, we can do
better.  This section gives a separate Boolean-gate attack for one certified
official Level-I public key.  It uses the May--Ostuzzi--Ressler transformation
only as a target-independent public change of coordinates; the analysis does
\emph{not} use their Assumption~1 or a random-tree model.

Write the quotient map in its two Frobenius channels as
\[
  D:A^{64}\longrightarrow A^{16},\qquad
  H:A^{64}\longrightarrow\F_{19}^{16},
  \qquad A=\F_{19^2},\quad Q=|A|=361.
\]
The retained target is $(t_D,t_H)$.  As in the rest of the paper, the target
sampler is rejection-corrected so that the retained coordinates are exactly
uniform.  The two channel families are different linear images of the same
native byte-reduced public coefficients and are therefore not independent.
Our analysis conditions first on the complete diagonal channel.  This fixes
the MOR staircase, the JG slice, and every diagonal certificate.  Only then do
we expose the residual Hermitian channel block by block.  This order is
important: it makes the key dependence of the slice harmless and permits the
conditional joint-pencil theorem below.

\subsection{The 15-form redesign and its fixed-key rank certificate}

Instead of transforming all sixteen diagonal equations, omit one original
base form and apply Just Guess to the other fifteen with
\begin{equation}
  m'=15,\qquad k=4,\qquad p=7,\qquad m'-k-p=4.
  \label{eq:jg-redesign-parameters}
\end{equation}
For the certified key used here, the published MOR transformation succeeds
with these parameters.  Extend its invertible $15\times15$ output change by
the untouched omitted form.  The resulting sixteen target coordinates are
again uniform and independent because this is an invertible public linear
change of the original diagonal target.

The transformed solver streams nine $A$-coordinates, guesses four, follows
all roots of seven univariate quadratics, enumerates the complete affine
solution space of the residual $4\times4$ linear system, checks the four
remaining transformed equations, and finally checks the omitted original
form.  Thus the outer population remains
\begin{equation}
  Q^{9+4}=Q^{13},                                  \label{eq:jg-outer-pop}
\end{equation}
exactly as in the earlier $(16,5,6)$ route.

\paragraph{Why the four guesses are not an arbitrary choice.}
The artifact also performs an exact parameter sweep for this fixed key under
the published deterministic MOR construction and its stated form order.  For
every one of the sixteen choices of an omitted diagonal base form, every
admissible fifteen-form choice with $k=3$ and
$p\in\{9,10,11,12\}$ fails.  For $p=9$, all sixteen omissions reach column~8
with a $57\times64$ constraint matrix of rank~$56$ and an eight-dimensional
kernel; for the certified omission, the same rank-$56$/kernel-$8$ obstruction
persists for $p=9,10,11,12$.  The full sixteen-form system likewise fails for
every admissible deterministic choice with $k\le4$; its first successful
setting is the known $(k,p)=(5,6)$ route.  Thus the exponent $Q^{13}$ is not
merely the first parameter choice found by a local search.  We do not promote
this finite sweep to an impossibility theorem for arbitrary reorderings or
arbitrary preliminary output recombinations.

Let $L_{\rm JG}\cong A^{24}$ be the direction space consisting of the fifteen
MOR coordinates together with the nine streamed coordinates.  A trial fixes
the other forty $A$-coordinates uniformly, hence samples a uniform affine
$L_{\rm JG}$-coset.  For $\lambda\in A^{16}$ let
\[
  K_{\rm full}(\lambda)\in A^{64\times24}
\]
be the polar matrix of the scalar diagonal combination
$\lambda\cdot D$, with the second argument restricted to $L_{\rm JG}$ and the
first argument left in the full $A^{64}$ ambient space.

\begin{proposition}[Certified maximum-rank diagonal pencil]
\label{prop:jg-rank24}
For the official Level-I $\ell=2$ KAT public key with count~$0$ and the
explicit $(15,4,7)$ MOR staircase specified in the artifact,
\begin{equation}
  \boxed{\rank_A K_{\rm full}(\lambda)=24
         \quad\text{for every }0\ne\lambda\in A^{16}.}
  \label{eq:jg-rank24}
\end{equation}
This is the maximum possible rank because $K_{\rm full}(\lambda)$ has only
$24$ columns.  The certificate uses exact finite-field arithmetic; no random
sampling enters the claim.  It is a fixed-key statement, not a theorem about
the probability that a fresh key has the same property.
\end{proposition}

\begin{proof}[Certificate summary]
We reconstruct the count-$0$ public key directly from the serialized official
KAT record and derive its sixteen $A$-valued diagonal quadratics in physical
coordinates.  The channel basis is
\[
 u\longleftrightarrow
 \begin{pmatrix}6&1\\1&13\end{pmatrix},\qquad
 L_D=\begin{pmatrix}11&5&4\\15&7&16\end{pmatrix},
\]
and direct evaluation agrees coefficient-for-coefficient with the
$A$-valued system used by the MOR certificate.  We then omit original form~0,
run the deterministic MOR transformation at $(15,4,7)$, and adjoin the
omitted form after the same domain transformation.

Fix a projective chart
$e=\min\operatorname{supp}(\lambda)$ and normalize $\lambda_e=1$; put
$s=15-e$.  To rule out rank below $24$, introduce a projective
$v\in A^{24}$ and the $64$ bilinear equations
$K_{\rm full}(\lambda)v=0$.  Multiplying the constant-chart equations by the
$24$ coordinates of $v$ gives
\[
 A_0:A^{64}\otimes A^{24}\longrightarrow \operatorname{Sym}^2(A^{24}),
 \qquad 1536\longrightarrow300.
\]
On every chart the exact rank of $A_0$ is $300$, leaving a kernel $N$ of
dimension $1236$.  For each free projective coordinate $i$ let
$A_i$ be the corresponding multiplication map and put
$C_i=A_iN\in A^{300\times1236}$.

We use a bidegree-$(s,2)$ block-Macaulay certificate.  Its only nontrivial
combinatorial ingredient is a capacity-two orientation from monomials of
$\lambda$-degree $s$ to predecessors of degree $s-1$.  To see that such an
orientation exists, consider any set $B$ of degree-$s$ monomials and count
predecessor edges with multiplicity.  Every monomial in $B$ has total
multiplicity $s$, while a degree-$(s-1)$ predecessor receives total
multiplicity at most $(s-1)+s=2s-1$.  Hence
\[
   s|B|\le(2s-1)|N(B)|<2s|N(B)|,
\]
so $|B|\le2|N(B)|$.  Hall's theorem with capacity two gives the required
orientation.  (For $s=0$ the chart is constant.)

After pivoting the same-degree target blocks with $A_0$, a degree-$(s-1)$
source monomial is therefore responsible for at most two top-degree target
blocks.  Thus top-degree surjectivity reduces to the finite assertions
\[
 \rank C_i=300,
 \qquad
 \rank\!\begin{bmatrix}C_i\\C_j\end{bmatrix}=600
 \quad(i\ne j).
\]
The artifact checks every such assertion on every chart: $120$ single-block
and $560$ pair-block checks, $680$ exact ranks in total, with no failures.
The largest rank matrix is only $600\times1236$ over $A$.  Block-triangular
elimination then makes the complete Macaulay target surjective.  In
particular every quadratic monomial in $v$ lies in the generated ideal, so a
projective common zero would force $v=0$, a contradiction.  Hence
\eqref{eq:jg-rank24} holds on every chart.
\end{proof}

\begin{table}[t]
\centering
\caption{Exact fixed-key certificate for \eqref{eq:jg-rank24}.  On chart
$e$, $s=15-e$ free projective coordinates give $s$ single-block and
$\binom{s}{2}$ pair-block rank checks.  All $680$ checks have the displayed
full row rank.}
\label{tab:jg-rank24-certificate}
\small
\begin{tabular}{@{}l r r r@{}}
\toprule
Certificate item & Matrix size & Number & Failures\\
\midrule
Constant multiplication maps $A_0$ & $300\times1536$ & $16$ & $0$\\
Single $C_i$ blocks & $300\times1236$ & $120$ & $0$\\
Pair stacks $[C_i;C_j]$ & $600\times1236$ & $560$ & $0$\\
\midrule
Finite rank checks & --- & $680$ & $0$\\
\bottomrule
\end{tabular}
\end{table}

\paragraph{Official-key anchoring and a 100-key staircase census.}
The certified public instance is count~$0$ in
\texttt{PQCsignKAT\_SNOVA\_48\_16\_19\_2.rsp}.  The artifact records SHA-256
\texttt{cb4951074a1a28366617cf3cc2f3d64572dff2daab89045f5bf0f9606f3c5627}
for that KAT file and
\texttt{ce99d967cfdc849441981be59fe50cbb92c4faacf2b523e0526f62c470ed9d7a}
for the count-$0$ serialized public key.  Independently of the rank-$24$
certificate, we expanded all $100$ public keys in this official KAT corpus and
reran the same deterministic fifteen-form MOR construction from their public
matrices.  All $100/100$ admit the $(15,4,7)$ staircase.  This finite census
is evidence against a rare-staircase explanation; it is not promoted to a
fresh-key probability theorem, and the rank-$24$ certificate itself is
claimed only for count~$0$.

\paragraph{Cost and space of the public preflight.}
The certificate is not treated as a free oracle.  A deliberately naive
elimination ledger already suffices: checking all $560$ pair stacks by dense
Gaussian elimination on at most $600\times1236$ matrices, together with the
$A_0$ and single-block checks and charging fewer than $2^{10}$ AXN gates per
$A$-field operation, is below
\begin{equation}
  C_{\rm pre}<2^{50}.                                \label{eq:jg-preflight-cost}
\end{equation}
Unlike the former rank-$17$ certificate, this proof requires no
$17017\times17108$ termination matrix.  A verifier may regenerate one chart
and one pair at a time.  With 16-bit storage for an $A$ symbol, the largest
kernel basis has $1536\cdot1236$ entries and the largest rank stack has only
$600\cdot1236$ entries; allowing separate input, elimination, and temporary
buffers keeps a straightforward implementation below a conservative $64$ MB
of mutable certificate-verification state.  This is a bound on verifier
working storage, not a claim about cache behavior or wall-clock time.

The maximum-rank certificate gives an essentially optimal diagonal second
moment.  For $0\ne y\in L_{\rm JG}$ let
\[
  T_y:A^{64}\to A^{16}
\]
be the polar derivative of $D$ in direction $y$.  Double-counting annihilating
output functionals gives the exact dual identity
\begin{equation}
  1+\sum_{0\ne y\in L_{\rm JG}}Q^{-\rank_A T_y}
  =Q^{-16}\sum_{\lambda\in A^{16}}
       Q^{24-\rank_A K_{\rm full}(\lambda)}.         \label{eq:jg-dual-rank-identity}
\end{equation}
Because every nonzero pencil element has rank exactly $24$,
\begin{equation}
 \eta_D:=\sum_{0\ne y\in L_{\rm JG}}Q^{-\rank_A T_y}
 = Q^8-Q^{-16}.                                      \label{eq:jg-eta}
\end{equation}
Apply \Cref{thm:collision-identity} over $A$ with ambient space $A^{64}$,
slice direction $A^{24}$, and sixteen diagonal outputs.  If $Z$ is the number
of full diagonal roots in a random trial and $\mu=Q^8$, then
\begin{equation}
 \mathbb EZ=\mu,\qquad
 \mathbb EZ^2=Q^{16}+Q^8-Q^{-8},\qquad
 \frac{\operatorname{Var}Z}{\mu^2}
 =\delta_{24}:=Q^{-8}-Q^{-24}.                       \label{eq:jg-variance}
\end{equation}
Thus the old $Q^{-1}$-scale variance allowance disappears entirely.

\subsection{Exact aggregate Just-Guess counting}

\begin{lemma}[Aggregate target-averaged Just-Guess count]
\label{lem:jg-aggregate-count}
Fix the public diagonal system, the certified target-independent
$(15,4,7)$ staircase, and the affine values of the forty non-streamed
coordinates.  Across all $Q^9$ streamed assignments and all $Q^4$ guesses,
let $N_j$ be the number of live nodes after $j$ univariate quadratic stages.
If every root is followed, including all $Q$ roots of an identically zero
equation, then
\begin{equation}
  \mathbb E N_j=Q^{13}\qquad(0\le j\le7).            \label{eq:jg-node-count}
\end{equation}
If every solution of every residual four-variable linear system is enumerated,
the expected total number $R$ of completions is
\begin{equation}
  \mathbb E R=Q^{13}.                                \label{eq:jg-linear-count}
\end{equation}
After the four transformed checks and the omitted original diagonal equation,
\begin{equation}
  \mathbb E Z=Q^8=19^{16}.                           \label{eq:jg-Droot-count}
\end{equation}
No independence between branches and no random-MQ hypothesis is used.
\end{lemma}

\begin{proof}
A prospective depth-$j$ node is specified by the thirteen outer coordinates
(nine streamed coordinates and four guesses) and the $j$ already solved
quadratic coordinates, hence by a point of $A^{13+j}$.  For a fixed such
point, the first $j$ transformed equations prescribe $j$ coordinates of the
uniform transformed target, so its probability of being live is $Q^{-j}$.
Linearity of expectation proves \eqref{eq:jg-node-count}.

A residual completion is specified by all twenty-four slice coordinates and
satisfies the first $7+4=11$ transformed selected equations.  Each fixed point
therefore contributes with probability $Q^{-11}$, giving
$Q^{24}Q^{-11}=Q^{13}$.  Requiring the remaining four transformed coordinates
and the omitted original coordinate prescribes all sixteen diagonal target
coordinates and gives $Q^{24}Q^{-16}=Q^8$.
\end{proof}

\subsection{A conditional full-rank theorem for the joint $D/H$ pencil}
\label{sec:jg-joint-pencil}

The previous pairwise filtering analysis bounded collisions after conditioning
on the diagonal channel.  A stronger argument works at the level of the
entire joint polar pencil and removes the pairwise cross-channel second-moment
loss.

Condition on the complete native diagonal-channel coefficient family $D$.
The deterministic MOR staircase, its domain transform, the direction space
$L_{\rm JG}$, and the rank-$24$ certificate of
\Cref{prop:jg-rank24} are then fixed.  The residual Hermitian coefficients of
different native blocks and different base forms remain independent in the
pinned XOF-transcript model.

\begin{lemma}[Conditional Hermitian block atoms]
\label{lem:jg-block-atoms}
Let
\[
  \rho_H=\frac{49}{829}<0.059108.
\]
Condition on the complete diagonal coefficient of one native block.  For a
diagonal block, every atom of its remaining one-dimensional Hermitian
coefficient has probability at most $\rho_H$.  For an off-diagonal block,
every nonzero scalar functional of its remaining two-dimensional Hermitian
coefficient has atom at most $\rho_H$, while every full two-dimensional point
has conditional probability at most
\begin{equation}
  \sigma_H=\frac{38416}{11886083}
  <\left(\frac{49}{829}\right)^2=\rho_H^2.
  \label{eq:jg-joint-block-atom}
\end{equation}
Conditioned on $D$, the residual blocks are independent across native blocks
and base forms.  Consequently, if an affine linear system of constraints in
these residual Hermitian blocks has base-field rank $r$, its conditional
probability is at most $\rho_H^r$.
\end{lemma}

\begin{proof}
In the diagonalizing coordinates used throughout the artifact, a diagonal
native block satisfies
\[
 D=(a+7b+16c)+(2b+7c)u,\qquad H=a+7b+18c,
\]
while an off-diagonal block satisfies
\[
 D=(a+13b+13c+16d)+(b+c+7d)u,
\]
\[
 H=(a+13b+13c+18d)+(18b+c)u.
\]
A raw residue has weight $14$ for $0,\ldots,8$ and $13$ for
$9,\ldots,18$.  Exact weighted enumeration of all $19^3$ diagonal tuples
gives maximum conditional atom
$2744/46424=49/829$.  Exact weighted enumeration of all $19^4$
off-diagonal tuples gives maximum conditional point mass
$38416/11886083$, and enumeration of the twenty projective nonzero scalar
functionals gives a smaller one-dimensional maximum than $49/829$.
The displayed inequality in \eqref{eq:jg-joint-block-atom} is exact.

For the last statement, expose independent native blocks one at a time and
perform block Gaussian elimination on the affine system.  A block can
increase the constraint rank by $0$, $1$, or (for an off-diagonal block) $2$.
A rank-one increment forces one nonzero scalar functional and costs at most
$\rho_H$; a rank-two increment fixes the entire two-dimensional block and
costs at most $\sigma_H<\rho_H^2$.  Multiplication over the independent blocks
gives $\rho_H^r$.
\end{proof}

\begin{lemma}[Cross-tensor injection]
\label{lem:jg-cross-tensor}
Let $U$ be an $A$-vector space, $0\ne y\in U$, and
$0\ne c\in\operatorname{Sym}^2_{\F_{19}}(A)$.  Writing
$c=\sum_\nu a_\nu\mathbin\odot b_\nu$, define
\[
 \Xi_{c,y}(x)=\sum_\nu\bigl((a_\nu y)\mathbin\odot(b_\nu x)
 +(b_\nu y)\mathbin\odot(a_\nu x)\bigr)
 \in\operatorname{Sym}^2_{\F_{19}}(U).
\]
Then
\begin{equation}
  \ker\Xi_{c,y}\subseteq Ay.
  \label{eq:jg-cross-tensor-kernel}
\end{equation}
In particular, for $U=A^{64}$ the base-field rank of $\Xi_{c,y}$ is at least
$128-2=126$.
\end{lemma}

\begin{proof}
If $x\notin Ay$, project $\operatorname{Sym}^2(U)$ onto the cross summand
between the disjoint $A$-lines $Ay$ and $Ax$ and identify that summand with
$Ay\otimes Ax$.  The projection of $\Xi_{c,y}(x)$ is the ordinary
symmetrization of $c$.  In odd characteristic this symmetrization is
injective on $\operatorname{Sym}^2(A)$: diagonal tensors acquire the nonzero
factor~$2$, while the off-diagonal symmetric tensors remain distinct.
Therefore the projection is nonzero whenever $c\ne0$.
\end{proof}

For a base-field output functional
$\lambda=(\lambda_D,\lambda_H)\in\F_{19}^{32}\oplus\F_{19}^{16}$, let
\[
 K_{D,H}(\lambda)\in\F_{19}^{128\times48}
\]
be the polar matrix of the scalar joint combination
$\lambda_D\cdot D+\lambda_H\cdot H$, with its second argument restricted to
$L_{\rm JG}$ and its first argument left in the full $128$-dimensional
ambient space.

\begin{theorem}[Conditional full-rank joint pencil]
\label{thm:jg-joint-fullrank}
Fix an arbitrary diagonal transcript $D$ for which the supplied
$(15,4,7)$ staircase and \Cref{prop:jg-rank24} hold.  In the conditional
native-coordinate XOF-transcript model for the remaining Hermitian channel,
\begin{equation}
 \Pr\!\left[\exists\,0\ne\lambda\in\F_{19}^{48},\ 0\ne y\in L_{\rm JG}:
     K_{D,H}(\lambda)y=0\ \middle|\ D\right]
 <2^{-114.684}.
 \label{eq:jg-joint-badkey}
\end{equation}
Equivalently, outside this event,
\begin{equation}
  \boxed{\rank_{\F_{19}}K_{D,H}(\lambda)=48
         \quad\text{for every }0\ne\lambda\in\F_{19}^{48}.}
  \label{eq:jg-joint-rank48}
\end{equation}
The bound is uniform in the conditioned diagonal transcript; no independence
between the diagonal coefficients and the key-dependent JG slice is assumed.
\end{theorem}

\begin{proof}
First suppose $\lambda_H=0$.  Every base-field functional on $A^{16}$ has the
form $z\mapsto\operatorname{Tr}_{A/\F_{19}}(\alpha\cdot z)$ for a unique
$\alpha\in A^{16}$.  The trace pairing on $A$ is nondegenerate, so the
rank-$24$ $A$-linear statement in \Cref{prop:jg-rank24} implies base-field
rank~$48$ for every nonzero diagonal-only functional.

Now fix a projective pair $([\lambda],[y])$ with $\lambda_H\ne0$.  Choose a
base form $i$ with nonzero Hermitian coefficient in $\lambda_H$, and condition
further on every residual Hermitian block except those of this one base form.
For each ambient $x\in A^{64}$, the coefficient vector of the remaining
Hermitian symbols in the scalar polar equation
$\lambda\cdot B_{D,H}(x,y)=0$ is exactly the cross tensor
$\Xi_{c,y}(x)$ for a nonzero
$c\in\operatorname{Sym}^2_{\F_{19}}(A)$.  By
\Cref{lem:jg-cross-tensor}, this coefficient map has base-field rank at least
$126$.  Hence the event $K_{D,H}(\lambda)y=0$ forces a rank-$126$ affine
system in the independent residual blocks of form~$i$.  By
\Cref{lem:jg-block-atoms}, its conditional probability is at most
$\rho_H^{126}$.

The number of projective points in $\F_{19}^{48}$ is
\[
  P_{47}=\frac{19^{48}-1}{18}.
\]
Conditioning on $D$ fixes $L_{\rm JG}$, whose base-field dimension is $48$,
so there are at most $P_{47}^2$ projective pairs to union-bound.  Therefore
\[
 \Pr[\text{bad joint pencil}\mid D]
 \le P_{47}^2\left(\frac{49}{829}\right)^{126}
 <2^{-114.684},
\]
which proves \eqref{eq:jg-joint-badkey}.  The bilinearity of the polar tensor
makes projectivization valid independently in $\lambda$ and $y$.
\end{proof}

\begin{corollary}[Exact joint-root second moment]
\label{cor:jg-joint-second-moment}
Fix a public key satisfying \eqref{eq:jg-joint-rank48}.  Let $C$ be a uniform
affine $L_{\rm JG}$-coset and let $t=(t_D,t_H)$ be the independent uniform
retained verifier target.  If
\[
 Z_{D,H}=|\{x\in C:(D(x),H(x))=t\}|,
\]
then
\begin{equation}
  \mathbb E Z_{D,H}=1,\qquad
  \mathbb E Z_{D,H}^2=2-19^{-48},
  \label{eq:jg-joint-exact-moments}
\end{equation}
and consequently
\begin{equation}
  \boxed{\Pr[Z_{D,H}>0]\ge\frac{1}{2-19^{-48}}>\frac12.}
  \label{eq:jg-joint-half-success}
\end{equation}
\end{corollary}

\begin{proof}
The first moment is $19^{48}19^{-48}=1$.  For the second moment, average over
a uniform affine coset and then over all ordered pairs in that coset.  A pair
has difference $y\in L_{\rm JG}$.  The diagonal $y=0$ contributes~$1$.  For
each $0\ne y\in L_{\rm JG}$, choosing the first point uniformly over the full
ambient space makes the collision equation
\[
 (D(x+y),H(x+y))=(D(x),H(x))
\]
an affine system whose linear part is the joint derivative in direction $y$.
Equation \eqref{eq:jg-joint-rank48} (equivalently its dual no-zero-divisor
form) gives rank~$48$, so the collision probability is $19^{-48}$.  The
independent uniform target contributes another factor $19^{-48}$.  Summing
the $19^{48}(19^{48}-1)$ ordered off-diagonal pairs gives
$1-19^{-48}$, proving \eqref{eq:jg-joint-exact-moments}.  Cauchy--Schwarz
(Paley--Zygmund at threshold zero) gives
$\Pr[Z_{D,H}>0]\ge(\mathbb EZ_{D,H})^2/\mathbb EZ_{D,H}^2$.
\end{proof}

\paragraph{What remains modeled.}
The theorem conditions on the entire diagonal transcript before exposing the
Hermitian residuals, so the key-dependent staircase and slice create no
independence gap.  It is nevertheless a distributional statement about those
residual XOF symbols, not a deterministic certificate that the already fixed
count-$0$ Hermitian channel satisfies \eqref{eq:jg-joint-rank48}.  The
artifact separately reconstructs that actual channel and includes
reconnaissance rank tests, but no sampled test is used as a theorem.  The
fixed-key diagonal rank-$24$ result and the $100/100$ staircase census remain
exactly as stated above.

\begin{proposition}[Depth-first space bound for the concrete Just-Guess route]
\label{prop:jg-lowspace}
The online $(15,4,7)$ enumeration can be implemented without storing any
$Q^{13}$-scale candidate list.  Excluding the immutable public key and the
separately verified public certificate, fewer than $2^{10}$ live $A$-symbols
and fewer than $2^{16}$ control bits suffice for the solver traversal.
Consequently the online attack has polynomial, in fact tiny, mutable state;
its exponential resource is time, not space.
\end{proposition}

\begin{proof}
Enumerate the nine streamed coordinates, four guessed coordinates, and at
most four residual-kernel parameters with nested base-$361$ counters.  At a
quadratic stage keep only the current branch and a root iterator.  A
nondegenerate quadratic has at most two roots; an identically zero quadratic
is traversed by a base-$361$ counter rather than materializing $361$ children.
The residual solver needs only its $4\times5$ augmented matrix, one particular
solution, and at most four kernel vectors.  The current point contributes
$64$ $A$-symbols, and the seven-stage branch stack stores only a constant
number of coefficients/counters per level.  Even retaining separate input,
RREF, target, reconstruction, and branch buffers stays far below $2^{10}$
$A$-symbols.  The explicit loop/branch state used in the Boolean ledger is
also below $2^{16}$ bits.  Completed candidates are checked and discarded
immediately.  No step requires a table proportional to the number of streamed
assignments, guesses, or diagonal roots.
\end{proof}

\subsection{Finite Boolean-gate instantiation}

We use the constructive $692$-gate multiplication recurrence and $84$-gate
adder count over $A$, together with the verified monic quadratic-root circuit
\begin{equation}
  C_{\rm root}=16{,}886\quad\AXN\text{ gates}.       \label{eq:jg-root-gates}
\end{equation}
The same MOR arithmetic formulas that regenerate the earlier $(16,5,6)$
ledger give, for $(15,4,7)$, $356{,}034$ gates for the seven quadratic
substitutions/root calls, $114{,}688$ gates for emitted-child generation, and
$1{,}089{,}920$ gates to recover the four residual affine equations.

The redesigned residual has only four unknowns, so we replace the inherited
$5\times5$ singular-system reserve by an explicit fixed-schedule circuit.

\begin{lemma}[Fixed-schedule $4\times4$ residual solver]
\label{lem:jg-rref4}
A complete rank-deficient-safe reduction of a $4\times4$ affine system over
$A$, including consistency testing, one particular solution, and an affine
kernel basis, can be implemented with fewer than $2^{17}$ AXN gates.
\end{lemma}

\begin{proof}
Store the augmented matrix as four rows of five $A$-symbols.  Use the safe
inverse $x\mapsto x^{359}$, which is $x^{-1}$ for $x\ne0$ and is zero at
$x=0$.  One explicit thirteen-multiplication addition chain has exponents
$2,3,6,7,14,28,56,112,224,336,350,357,359$.  A four-column fixed schedule
therefore needs at most
\[
 4\cdot13+4\cdot5+4\cdot3\cdot5=132
\]
$A$-multiplications and $4\cdot3\cdot5=60$ $A$-additions for inversion,
normalization, and elimination.  This costs
\[
 132\cdot692+60\cdot84=96{,}384\quad\AXN\text{ gates}.
\]
For completeness, encode an $A$-symbol by its two five-bit $\F_{19}$
coordinates.  A ten-bit zero test costs at most $32$ AXN gates, a conditional
select of an $A$-symbol at most $30$, and a conditional swap at most $40$.
Scanning at most four candidate rows in each of four columns costs at most
$16\cdot32$ gates for zero tests.  To avoid any hidden dynamic-addressing
assumption, pessimistically implement all four possible pivot-slot swaps and
mask the unused ones; together with no-pivot masking, consistency logic, and
extraction of one particular solution plus up to four kernel vectors, the
complete fixed schedule remains below $124{,}384<2^{17}$ AXN gates.  Singular
systems require no exceptional branch.
\end{proof}

The remaining address and dispatch logic is likewise small once the attack is
implemented depth first.  The nine stream coordinates, four guesses, and at
most four kernel parameters are nested base-$361$ counters; emitted-child
storage is already charged in the child-generation line, and pivot control is
already charged in \Cref{lem:jg-rref4}.  Using $128$ gates per update of any
of the seventeen base-$361$ loop digits gives $2176$ gates.  We charge $2048$
gates at each of the seven quadratic stages for branch/stack dispatch, $3584$
gates for all five sequential-check guards and restart flags, and $4096$ gates
for residual/kernel-enumeration dispatch.  The explicit subtotal is therefore
\[
 2176+7\cdot2048+3584+4096=24{,}192<2^{15}.
\]
We nevertheless reserve
\begin{equation}
 C_{\rm ctrl}=2^{17}                                \label{eq:jg-control-reserve}
\end{equation}
to leave more than a factor-four margin over that itemization.

For the final diagonal checks, one transformed selected equation costs at
most $104{,}676$ gates under the same MOR formula.  The omitted original form
need not inherit the staircase, so we instead charge the generic dense
$A^{24}$ evaluation cap
\begin{equation}
 C_{\rm eval}\le459{,}024<2^{19}.                    \label{eq:jg-eval-gates}
\end{equation}
Evaluate the four cheaper transformed checks first and the omitted form last.
By the same target-averaging argument as \Cref{lem:jg-aggregate-count}, the
expected check work per expected $Q^{13}$ completion population is bounded by
\begin{equation}
 104{,}676(1+Q^{-1}+Q^{-2}+Q^{-3})
   +459{,}024Q^{-4}
 <104{,}967.                                         \label{eq:jg-sequential-checks}
\end{equation}
No independence between candidates is used; a point reaches check $j$ exactly
when the preceding target coordinates match.

\begin{table}[t]
\centering
\caption{Audited AXN ledger per expected $Q^{13}$ aggregate population for
the certified $(15,4,7)$ route.}
\label{tab:jg-redesign-ledger}
\small
\begin{tabular}{@{}l r@{}}
\toprule
Component & AXN gates\\
\midrule
Seven quadratic substitutions and root circuits & $356{,}034$\\
Seven child-generation reserves & $114{,}688$\\
Recover four residual affine equations & $1{,}089{,}920$\\
Fixed-schedule $4\times4$ RREF/kernel reserve & $131{,}072$\\
Completion generation & $19{,}400$\\
Sequential four selected checks plus omitted form & $104{,}967$\\
Remaining control/address/dispatch reserve & $131{,}072$\\
\midrule
Total & $\mathbf{1{,}947{,}153}<2^{21}$\\
\bottomrule
\end{tabular}
\end{table}

We retain a $2^{40}$-gate cap per surviving diagonal root for $H$, formatting,
reconstruction, hashing, and the original verifier.  We also allow the looser
$Q^9 2^{40}$ reserve for chosen-message target generation.  These terms are
tiny compared with $Q^{13}2^{21}$ but are included below.

\begin{proposition}[Certified-diagonal Level-I Boolean bound]
\label{prop:jg-level1-gates}
For a Level-I diagonal transcript satisfying \Cref{prop:jg-rank24}, and for
any Hermitian completion satisfying \eqref{eq:jg-joint-rank48}, one complete
target-averaged $(15,4,7)$
trial has expected work
\begin{equation}
 \mathbb E W_{\rm JG}^{\rm trial}
 <Q^{13}2^{21}+Q^8 2^{40}+Q^9 2^{40}.                \label{eq:jg-expected-work}
\end{equation}
The algorithm follows all quadratic branches and enumerates every consistent
residual linear completion; May--Ostuzzi--Ressler Assumption~1 is not used.
By \Cref{cor:jg-joint-second-moment}, the per-trial success probability is
at least $(2-19^{-48})^{-1}>1/2$.  Hence
\begin{equation}
 \boxed{\log_2\mathbb E W_{\rm JG}^{\rm success}<132.447.}
 \label{eq:jg-numerical}
\end{equation}
Thus the clean bound is more than $10.55$ bits below the Level-I classical
reference $143$.  Under the conditional Hermitian-transcript distribution of
\Cref{thm:jg-joint-fullrank}, the exceptional public-key event has probability
below $2^{-114.684}$.  Using the literal component total $1{,}947{,}153$ instead
of rounding it to $2^{21}$ gives
\begin{equation}
  \log_2\mathbb E W_{\rm JG}^{\rm success}<132.340,
  \label{eq:jg-literal-numerical}
\end{equation}
about $10.66$ bits below the reference.  The one-time public-certificate check
costs less than $2^{50}$ gates by \eqref{eq:jg-preflight-cost} and is therefore
negligible on this scale even when charged explicitly.  The proof applies verbatim to any other Level-I diagonal transcript on which
the same deterministic staircase and rank-$24$ predicates pass, for every
Hermitian completion satisfying \eqref{eq:jg-joint-rank48}.  The conditional
model bounds failure of that latter predicate by \eqref{eq:jg-joint-badkey};
no probability that a fresh key passes the diagonal certificate is asserted.
\end{proposition}

\begin{proof}
The aggregate work through all five diagonal checks is at most
$1{,}947{,}153Q^{13}$ by
\Cref{lem:jg-aggregate-count,lem:jg-rref4,tab:jg-redesign-ledger}.  Adding the
two declared post-processing/target-generation reserves yields
\eqref{eq:jg-expected-work} after rounding the aggregate component allowance
upward to $2^{21}$.  For every Hermitian completion satisfying \eqref{eq:jg-joint-rank48}, the
per-trial success probability is at least the right side of
\eqref{eq:jg-joint-half-success}.

For independent trials, expected total cost to the first success is
$\mathbb EC_1/\Pr[S_1=1]$ even if within-trial cost and success are correlated:
$\{N\ge i\}$ depends only on earlier trials and is independent of the cost of
trial $i$.  Dividing \eqref{eq:jg-expected-work} by
$(2-19^{-48})^{-1}$ gives \eqref{eq:jg-numerical}; using the literal component
sum gives \eqref{eq:jg-literal-numerical}.  Finally add the one-time $2^{50}$
public preflight reserve; it changes neither displayed decimal exponent.
\end{proof}

\paragraph{Interpretation.}
The Level-I result now has neither a random-tree multiplier nor an unverified
candidate-abundance premise for the certified diagonal transcript.  The MOR
staircase and the full-diagonal rank-$24$ condition are finite public
certificates, and all within-diagonal candidate counts are target-averaged
identities.  The remaining distributional qualification is narrower: after
conditioning on that complete diagonal transcript, the joint $D/H$ full-rank
property fails with probability below $2^{-114.684}$ in the pinned
XOF-transcript model.  We do not claim that the already fixed count-$0$
Hermitian channel has a deterministic joint-pencil certificate, nor that a
random fresh key passes the rank-$24$ diagonal certificate with any particular
probability.  The online attack
needs only the deterministic certificate verification priced in
\eqref{eq:jg-preflight-cost}; we do not claim a generic procedure that finds an
alternative certificate when this fixed one fails.  Every final candidate is
still checked by the unmodified verifier.  The two $\ell=4$ Level-I rows
continue to retain $\kappa_{\rm hom}$ in their homotopy ledgers.

\section{Streamed-XL and low-state analyses}
\label{sec:low-state-xl}

The complete-square homotopy route supplies a conditional time bound under
$\mathsf H_{\rm hom}$, but it does not supply a low-memory implementation.
This section records a complementary
branch whose matrices can be streamed with linear vector state.  Its
correctness is exact once a finite quotient certificate is produced; its
numerical exponent retains a production selected-degree premise.  The two
branches therefore answer different questions rather than competing for one
headline.

\subsection{Macaulay matrices and the low-state branch}
XL turns quadratic equations into a larger linear system by multiplying them
by selected monomials.  A multigraded version tracks the degree in each
variable block and can produce a smaller matrix for block-structured systems
~\cite{Cabarcas2025SNOVA,FurueEtAl2025}.  Recovery does not require a
one-dimensional kernel.  Commuting multiplication matrices can also describe
all roots through FGLM and finite-field factorization
~\cite{FGLM1993,KehreinKreuzer2005}.  The analysis below states its additional
production-degree assumption.

\subsection{Eigenblocks and multidegrees}
Let $\Omega$ split the minimal polynomial of $S$, and let
$e_0,\ldots,e_{\ell-1}$ be the primitive idempotents of
$A\otimes_{\F_q}\Omega$.  If $L$ is $A$-linear of $A$-dimension $s$, then
\[
 L\otimes\Omega=L_0\oplus\cdots\oplus L_{\ell-1},
 \qquad \dim_\Omega L_a=s.
\]
The unordered power-pair features are an invertible recombination of the
unordered eigenblock features
\[
 x^\transpose\Sigma_{e_a}P_i\Sigma_{e_b}x,
 \qquad 0\le a\le b<\ell.
\]
They have multidegree $2\mathbf e_a$ when $a=b$ and
$\mathbf e_a+\mathbf e_b$ when $a<b$.  Translation creates lower-degree terms
only in the same one or two blocks.  One homogenizer per block therefore
preserves these exact multidegrees.  This is the grading used by the
streamed-XL branch.

\begin{proof}[Proof of the eigenblock statement]
The power basis and the idempotent basis of
$A\otimes\Omega\cong\Omega^\ell$ are related by an invertible Vandermonde
matrix.  Passing to the symmetric square gives an invertible change between
unordered power-pair and unordered eigenblock coordinates.  The idempotent
$e_a$ projects $L\otimes\Omega$ onto $L_a$.  A diagonal feature uses one
block twice; an off-diagonal feature is bilinear in two blocks.  Translation
and homogenization preserve the claimed supports.
\end{proof}

\subsection{Four-block special XL}
For an $A$-linear $\ell=4$ slice of $A$-dimension $s$, scalar extension gives
four blocks of $s$ variables.  Each unordered pair of blocks carries $m_1$
independent quadratic equations.  After adding one homogenizer per block, the
Hilbert series used by the published special-XL model is
\begin{equation}
 H_{m_1,s}(\mathbf t)=
 \frac{\prod_{a=1}^4(1-t_a^2)^{m_1}
       \prod_{1\le a<b\le4}(1-t_at_b)^{m_1}}
      {\prod_{a=1}^4(1-t_a)^{s+1}}.                 \label{eq:xl-hilbert}
\end{equation}
For a multidegree $\mathbf d=(d_1,\ldots,d_4)$, the monomial-column count is
\begin{equation}
  M_s(\mathbf d)=\prod_{a=1}^4\binom{s+d_a}{d_a}.    \label{eq:xl-columns}
\end{equation}
A negative coefficient of \eqref{eq:xl-hilbert} is the selected-degree trigger
in the cited model~\cite{Cabarcas2025SNOVA}.  It predicts that a planted
instance has the finite quotient needed for extraction; the theorem below
makes the operator and state accounting exact but does not prove that
production instances reach such a quotient at that degree.

The exact row count is
\begin{align}
 R_s(\mathbf d)
={}&m_1\sum_{a=1}^4
 \binom{s+d_a-2}{d_a-2}\prod_{c\ne a}\binom{s+d_c}{d_c}
 \notag\\
 &+m_1\sum_{1\le a<b\le4}
 \binom{s+d_a-1}{d_a-1}
 \binom{s+d_b-1}{d_b-1}
 \prod_{c\notin\{a,b\}}\binom{s+d_c}{d_c},          \label{eq:xl-rows}
\end{align}
where a binomial coefficient with a negative lower argument is zero.

\begin{proposition}[Streaming the special-XL operator]
\label{prop:streamed-xl}
Let $A_{\mathbf d}$ be the Macaulay matrix at multidegree $\mathbf d$, with
$M=M_s(\mathbf d)$ columns and $R=R_s(\mathbf d)$ rows.  Every row and
transpose-row product can be generated from the public quadratic coefficients
using $O(s^2)$ scratch space.

Over an extension $E/\F_{19^4}$, choose a random diagonal
$\Delta\in E^{R\times R}$ and put
\[
  B=A_{\mathbf d}^\transpose\Delta A_{\mathbf d}.
\]
Then
\[
  \Pr[\ker B\ne\ker A_{\mathbf d}]\le M/|E|.
\]
A product by $B$ is evaluated by one pass over the streamed rows and uses
$O(M)$ resident field-vector state.  Scalar Wiedemann over $E=\F_{19^8}$
requires fewer than
\begin{equation}
  16R(s+1)^2M                                      \label{eq:xl-stream-work}
\end{equation}
signed $\F_{19^4}$ multiply--accumulates per trial.  A conservative cubic
preconditioner over $\F_{19^{12}}$ gives the Level-I bound
$32R(s+1)^2M$.
\end{proposition}

\begin{proof}
The two sums in \eqref{eq:xl-rows} count multipliers of the diagonal and
off-diagonal equation families.  A homogenized quadratic row has at most
$(s+1)^2$ terms, so rows and transpose products are generated on demand.

Factor $A_{\mathbf d}=UV$ at rank $r$.  If
$U^\transpose\Delta U$ is nonsingular, then
$\ker(A_{\mathbf d}^\transpose\Delta A_{\mathbf d})=\ker V=\ker A_{\mathbf d}$.
By Cauchy--Binet, $\det(U^\transpose\Delta U)$ is a nonzero multilinear
polynomial of total degree $r$ in the diagonal entries.  Schwartz--Zippel
gives failure at most $r/|E|\le M/|E|$.

Streaming first computes the row inner product and then accumulates the
scaled row.  In the quadratic extension, the sparse dot product and
accumulation use at most four base-extension products per row term and the
diagonal multiplication uses three more.  Fewer than $3M$ normal-operator
products generate the scalar Wiedemann sequence and reconstruct one null
vector.  Rounding the resulting constant gives
\eqref{eq:xl-stream-work}; the cubic-extension calculation is analogous.
\end{proof}

\subsection{Certified extraction beyond exact corank one}
A one-dimensional Macaulay kernel is convenient but is not required for
soundness.

\begin{theorem}[Recovery from a finite commuting border quotient]
\label{thm:border-recovery}
Fix a target and slice over a finite field $\Omega$, and let
$I\subseteq\Omega[x_1,\ldots,x_h]$ be the ideal of the residual equations.
Let $\mathcal B$ be a finite divisor-closed monomial set containing $1$.
Suppose the Macaulay row space supplies, for every border monomial
$\beta\in\partial\mathcal B$, a certified relation
\[
  g_\beta=\beta-\sum_{b\in\mathcal B}c_{\beta,b}b\in I.
\]
Construct the formal multiplication matrices $M_1,\ldots,M_h$ on
$\Omega\mathcal B$.  If they commute, then the border relations generate a
zero-dimensional ideal $J\subseteq I$, $\mathcal B$ is a basis of
$\Omega[x]/J$, and every root of $I$ occurs among the roots enumerated from
the multiplication tables.  FGLM conversion costs $O(h|\mathcal B|^3)$ field
operations once the tables are available.  Complete residual and public
verification removes every extra point.
\end{theorem}

\begin{proof}
Commutativity makes $x_i\mapsto M_i$ an algebra homomorphism.  Divisor closure
identifies each basis monomial with its coordinate vector, and every border
relation acts as zero.  The border-division theorem therefore gives unique
normal forms in $\Span_\Omega\mathcal B$, so the relations form a border basis
of a zero-dimensional ideal $J$.  Because every relation is a Macaulay row
combination from $I$, we have $J\subseteq I$ and hence
$V(I)\subseteq V(J)$.  Standard FGLM conversion and finite-field
factorization enumerate the geometric points of $J$~\cite{FGLM1993,KehreinKreuzer2005}.
\end{proof}

A quotient of dimension larger than one is therefore a correctness fallback,
not a free extension of the scalar-Wiedemann exponent.  Obtaining all border
relations, block-nullspace work, and resident multiplication-table state must
fit the displayed ledger or be charged separately.

\subsection{Three-plus-one completion}
The fourth eigenblock can be completed linearly after solving the first three.
Let $E_{ab}$ be the $m_1$ equations supported on blocks $a,b$.

\begin{theorem}[Exact three-plus-one completion]
\label{thm:three-plus-one}
Suppose a certified first-stage solver returns a finite set containing every
root of the six families $E_{ab}$ with $1\le a\le b\le3$.  For each partial
root $p=(x_1,x_2,x_3)$, substitute into $E_{14},E_{24},E_{34}$ and obtain
\[
  A_px_4=b_p,
  \qquad A_p\in\Omega^{3m_1\times s}.
\]
Discard inconsistent systems.  If $\rank A_p=s$, solve for the unique fourth
block and retain it exactly when $E_{44}=0$.  If a consistent system has rank
below $s$, declare the trial incomplete unless a separately priced procedure
enumerates all completions.  Whenever no incomplete case occurs, the retained
set contains every root of the complete four-block system.  Final descent and
verification make the procedure Las Vegas.
\end{theorem}

\begin{proof}
Any complete root restricts to a first-stage root.  At that partial root its
true fourth block satisfies the displayed linear system.  On a trial that declares no incomplete case, every consistent system
encountered has full column rank.  Hence the true fourth block is the unique
completion, and the remaining diagonal family checks it.  Conversely every
retained point satisfies all ten families.
\end{proof}

\subsection{Status of the low-state branch}
The streamed operator, row count, and correctness fallbacks above remain
useful complementary machinery.  No production numerical profile is reported
because the artifact does not regenerate a selected multidegree,
success-probability bound, and state figure from one checker.  A numerical frontier would
require both that profile generator and a production border-basis or solve
transcript.


\section{A finite fixed-key spectrum-certificate format}\label{app:fixed-key-spectrum}
The idealized XOF-transcript theorem is not required for a statement about a concrete key.
The exact fixed-key rank-$24$ certificate used by the Level-I $(15,4,7)$
attack is summarized in \Cref{prop:jg-rank24} and supplied in the artifact.
This appendix gives the more general finite certificate format for aggregate
spectra.  The fixed-key statement concerns the certified public key only; no
fresh-key pass probability is claimed.

Put $N_U=\dim_{\F_q}U$.  Choose a basis
$y_0,\ldots,y_{h-1}$ of a fixed slice $L$ and write
\[
 T(Y)=\sum_{i=0}^{h-1}Y_iT_{y_i}
 \in\F_q[Y_0,\ldots,Y_{h-1}]^{K\times N_U}.
\]
For $0\le j<h$, use the disjoint normalized projective chart
$Y_0=\cdots=Y_{j-1}=0$, $Y_j=1$, and put
$R_j=\F_q[Y_{j+1},\ldots,Y_{h-1}]$.  Let $T^{(j)}$ be the resulting matrix.
For $t\ge0$, define
\begin{equation}
 I_{j,t}=\left\langle
   Y_i^q-Y_i\ (i>j),\quad
   \text{all }(t+1)\times(t+1)\text{ minors of }T^{(j)}
 \right\rangle\subseteq R_j.                        \label{eq:histogram-ideal}
\end{equation}
Put $n_{j,\le t}=\dim_{\F_q}(R_j/I_{j,t})$,
$n_{j,\le-1}=0$, and
\begin{equation}
  N_t=\sum_{j=0}^{h-1}(n_{j,\le t}-n_{j,\le t-1}). \label{eq:histogram-count}
\end{equation}

\begin{proposition}[Projective rank-histogram certificate]
\label{prop:projective-histogram}
The integer $N_t$ is exactly the number of projective directions
$[y]\in\mathbb P(L)(\F_q)$ for which $\rank T_y=t$.  Consequently
\begin{equation}
  \bar\eta_L^U=(q-1)\sum_tN_tq^{-t}.                 \label{eq:histogram-eta}
\end{equation}
Reduced Gr\"obner bases with their standard monomial sets, or polynomial
reduction transcripts proving the unit-ideal and quotient-dimension claims,
form machine-checkable fixed-key certificates.
\end{proposition}

\begin{proof}
The normalized charts contain exactly one representative of every projective
$\F_q$-direction.  The field equations identify
\[
 R_j/\langle Y_i^q-Y_i:i>j\rangle
 \cong\prod_{a\in\F_q^{h-j-1}}\F_q
\]
by evaluation.  Every ideal in this product of fields is radical, and its
quotient dimension is the number of coordinates on which all generators
vanish.  The minors in \eqref{eq:histogram-ideal} vanish exactly when
$\rank T^{(j)}(a)\le t$.  Thus $n_{j,\le t}$ counts rank-at-most-$t$ points
in chart $j$, and differencing gives \eqref{eq:histogram-count}.  Every
projective direction has $q-1$ nonzero scalar representatives and rank is
invariant under base-field scaling, proving \eqref{eq:histogram-eta}.
\end{proof}

A stronger certificate that $\rank T_y\ge t_0+1$ for every nonzero $y$
consists of unit-ideal transcripts for all $I_{j,t_0}$.  The complete histogram is more
flexible: a small exceptional low-rank locus can be charged exactly instead of
controlling the whole pencil by its minimum rank.

\section{Dimension formulas and exact numerical rows}\label{app:dimensions}
For a Version~2.3 row,
\[
 n=v+o,
 \quad m_1=\left\lceil\frac{or}{\ell}\right\rceil,
 \quad M=or\ell,
 \quad N_0=n\ell,
 \quad K=m_1\binom{\ell+1}{2}.
\]
For a full affine lift with $\rank C_{R,v}=M-K$,
\[
 \dim\ker C_{R,v}=N_0-M+K.
\]
For a structured lift whose complete offset orbit has rank $m_1\ell^2$,
\[
 \dim_AH^F_{R,v}=n-m_1\ell.
\]
The direct ordinary fresh-slice lower bound for $\ell=4$ is
\[
 \dim(V'\cap\ker C_{R,v})\ge4(v-1)-(M-K).
\]

\begin{table}[H]
\centering
\caption{Exact unit-leading exponents underlying the rounded main text.  Every
row conditions on its exact public quotient/rank preflights, the applicable
spectrum/model events, and $\mathsf H_{\rm hom}$; the $\ell=4$ rows additionally
condition on the fixed-layout structural preflight.  The values exclude
$\log_2\kappa_{\rm hom}$.  ``Direct'' solves the complete square; ``adaptive''
is the conditional mixture of an exact smaller-system determinant test and the
complete-square fallback.}
\label{tab:exact-ledger}
\small
\begin{tabular}{@{}c l r r r@{}}
\toprule
Level & Parameters & Direct & Adaptive & Adaptive headroom\\
\midrule
I   & $(28,5,4,4)$  & $142.606960$ & $128.246278$ & $14.753722$\\
I   & $(48,16,2,2)$ & $132.970670$ & $132.188072$ & $10.811928$\\
I   & $(28,4,4,5)$  & $142.606374$ & $128.245747$ & $14.754253$\\
III & $(40,7,4,4)$  & $185.450879$ & $171.616678$ & $35.383322$\\
III & $(72,24,2,2)$ & $183.794625$ & $183.794625$ & $23.205375$\\
III & $(38,5,4,5)$  & $185.450045$ & $171.615857$ & $35.384143$\\
V   & $(50,9,4,4)$  & $227.560406$ & $214.558321$ & $57.441679$\\
V   & $(96,32,2,2)$ & $233.843987$ & $233.843987$ & $38.156013$\\
V   & $(52,6,4,6)$  & $227.559359$ & $214.557277$ & $57.442723$\\
\bottomrule
\end{tabular}
\end{table}

All values charge the $150$, $692$, and $2628$-gate multiplication circuits
as appropriate and use the conservative $B^2$ forbidden-vector separator.
For $\ell=4$, the Jacobian-preflight density uses the projective right-kernel
union bound, not an independence claim after a biased linear transform.

\section{Arithmetic-circuit certificate details}\label{app:circuits}
The scalar multiplier consumes two little-endian five-bit canonical integers
in $\{0,\ldots,18\}$ and returns the canonical five-bit representation of
their product modulo $19$.  Its $150$ gates are two-input AND, XOR, and XNOR
gates with free constants, wires, fan-out, and structural sharing.  The
released netlist is tested on all $361$ valid input pairs.  The SHA-256 digest
of the canonical compact JSON payload containing exactly
\texttt{ninputs}, \texttt{gates}, and \texttt{outputs} is
\begin{center}
\ttfamily\small
 e1708a32c1475ad14295c2b34fed976fb09577b127810d29bf237db18aabc131
\end{center}
for the emitted balanced-netlist representation.  This is an internal payload
identifier, not the bytewise digest of the surrounding ledger file.

The tower recurrences are
\[
 \F_{19^2}=\F_{19}[u]/(u^2+1),
 \qquad
 \F_{19^4}=\F_{19^2}[v]/(v^2-(1+u)).
\]
Three-product Karatsuba at each level, together with the exact linear
operations stated in \Cref{thm:sq-circuits}, gives constructive recurrence
counts $692$ and $2628$.  The artifact verifies the tower identities, the
basis determinant, and the recurrence arithmetic; it does not emit or
exhaustively test complete extension-field netlists.

These circuits normalize only the pointwise multiplication schedule.  They
are not lower bounds and do not price the whole solver.  Additions,
inversions, factorization, bulk evaluation/interpolation transforms,
filtering, control,
indexing, communication, and memory remain in $\kappa_{\rm hom}$ and are
included in the unresolved multiplier $\kappa_{\rm hom}$; no numerical upper
bound on that multiplier is proved here.

\section{Reproduction and package contents}\label{app:artifact}
The source package is self-contained.  Its root contains the complete paper
source, bibliography, compiled PDF, build instructions, and hashes.  The
\texttt{artifact/} directory contains:
\begin{itemize}
\item a standalone generator and separately coded consistency checker for the
      all-nine conditional homotopy ledger;
\item the complete $150$-gate scalar multiplier netlist and tower checks;
\item the Just Guess expected-operation and capped ledgers, the quadratic-root
      check, and the exact conditional Frobenius-channel atom enumeration;
\item the official Level-I $\ell=2$ KAT count-$0$ reconstruction, the exact
      rank-$24$ diagonal-pencil certificate (all $680$ single/pair checks), and
      the $100/100$ deterministic $(15,4,7)$ staircase census;
\item a public-key correspondence harness for one official Level-I KAT; and
\item a non-planted zero-offset end-to-end composition test at the unofficial
      shape $(2,1,2,2)$, with fresh-key forgery and rank stress checks.
\end{itemize}

A clean reproduction is:
\begin{verbatim}
sha256sum -c SHA256SUMS
make regenerate
make verify
make
\end{verbatim}
The all-nine generator assigns no random probability to the fixed outer map:
its $\ell=4$ ledgers condition on the exact public structural preflight.  The
checkers are deterministic, require Python~3.10 or newer plus NumPy for the two
correspondence tests, and do not require network access.  The paper build uses
only the files in the package.  The supplied \texttt{SHA256SUMS} records every
distributed source and pinned input; generated TeX build files are excluded.

\end{document}
